% YAML frontmatter for the website — see template notes.
%--- iliad -------------------------------------------------------------------
% title: "Physics of Deep Learning"
% cluster: B
%--- end ----------------------------------------------------------------------
\documentclass[11pt]{article}

% iliad.sty: local copy first (Overleaf), else the shared ../iliad.sty
\IfFileExists{iliad.sty}{\usepackage[boxes]{iliad}}{\usepackage[boxes]{../iliad}}

% \solutionsfalse   % uncomment to hide all solutions from the PDF

% ---- Free zone -------------------------------------------------------------
\newcommand{\R}{\mathbb{R}}
\newcommand{\E}{\mathbb{E}}
\newcommand{\bv}[1]{\boldsymbol{#1}}
\newcommand{\ntk}{K_{\mathrm{NTK}}}

\title{Physics-Inspired Methods for Understanding Deep Learning}
\author{\authorname{Aleksander Czejdo} \and \authorname{Tudor Dimofte}\\ \affiliation{University of Edinburgh} \and \authorname{Brianna Grado-White} \and \authorname{Charles Renshaw-Whitman}} % TODO: fix author names/affiliations for this module
\date{\today}

\begin{document}
\maketitle

\begin{summary}
Physics-inspired methods for understanding deep learning, for PhD students
with math/physics backgrounds. The organizing claim, following the assigned
reading \cite{simon2026scientific}: deep learning exposes its own
``equations of motion,'' so the obstacle to a theory is complexity, not
opacity. These notes make three of that paper's five lines of evidence
rigorous --- solvable settings, insightful limits, and macroscopic laws ---
with emphasis on concrete payoffs: places where a physics-style calculation
predicts the behavior of real networks, in ways relevant to forecasting and
monitoring the capabilities of large models.
\end{summary}

\begin{learningoutcomes}
  \subsection*{Solvable settings (Block 1)}
  \begin{itemize}
    \item Compute learning curves for the teacher--student perceptron via the annealed approximation, and explain where the replica method is needed.
    \item Explain how deep linear networks decouple into solvable mode-wise ODEs, and why this produces a greedy low-rank bias.
    \item Derive the NTK description of a linearized network, and use the eigenlearning framework to trace the quantitative chain from architecture to inductive bias.
  \end{itemize}

  \subsection*{Insightful limits (Block 2)}
  \begin{itemize}
    \item Derive the order/chaos phase diagram of random deep networks and use it to predict trainable depth.
    \item Explain why depth-to-width ratio is the effective coupling of the finite-width expansion.
    \item Derive, heuristically, the $N\to\infty$ mean-field limit of one-hidden-layer networks, and explain how its time-dependent kernel encodes feature learning.
    \item Explain why intermediate scalings $\gamma \in [1/2,1)$ reduce to the kernel limit, and why layer-dependent learning-rate scalings ($\mu$P) are forced in deep networks --- enabling hyperparameter transfer.
  \end{itemize}

  \subsection*{Dynamics and macroscopic laws (Block 3)}
  \begin{itemize}
    \item Describe the DMFT order parameters of feature learning and what ``rich'' training does to them.
    \item Explain edge-of-stability dynamics and the central-flow resolution.
    \item Derive power-law scaling of test loss from power-law kernel spectra.
  \end{itemize}
\end{learningoutcomes}

%===============================================================================
\section{Prerequisites and how to use these notes}
\label{sec:prereqs}

Assumed background: statistical mechanics (free energies, saddle points,
phase transitions), probability, linear algebra; prior exposure to deep
linear networks and the decoupled ODEs of \cite{saxe2014exact} (we only
review their structure). Students read \S2 of \cite{simon2026scientific}
before the day begins. Phase transitions, singular learning theory, and
developmental interpretability are deferred to their own day, for which this
material is a prerequisite.

Suggested schedule for a one-day module:
\begin{itemize}
  \item 09:00--10:00 \quad Session 0: reading discussion, \cite{simon2026scientific} \S2 (\cref{sec:reading});
  \item 10:00--12:30 \quad Block 1: solvable settings (\cref{sec:block1});
  \item 12:30--13:30 \quad lunch; lab setup (\texttt{neural-tangents} installs);
  \item 13:30--14:30 \quad Lab: eigenlearning experiment (\cref{sec:lab});
  \item 14:30--16:15 \quad Block 2: limits, criticality, parameterizations (\cref{sec:block2});
  \item 16:30--18:00 \quad Block 3: dynamics and macroscopic laws (\cref{sec:block3});
  \item 18:00--18:45 \quad Closing: universality and alignment synthesis (\cref{sec:closing}).
\end{itemize}

\begin{callout}[note]
A word on notation. Different research communities studying deep learning
recycle the same letters with different meanings, and we have chosen to
keep each section close to the notation of its primary sources rather than
force a global convention; each section restates its conventions at the
start, and the following index collects the recurring symbols and the
known collisions.
\begin{itemize}
  \item $N$: the \emph{input dimension} in \cref{sec:statmech} (Block 1's
  perceptron sections) only; everywhere else, the \emph{width} of a hidden
  layer --- except \cref{sec:eft}, which follows its source and uses $n$
  for width and $L$ for depth (flagged there).
  \item Training-set size: $P$ in \cref{sec:statmech} (where $n$ is
  reserved for the replica index), $n$ in all kernel, lab, and scaling
  sections.
  \item $d$: input dimension (outside \cref{sec:statmech}); $\bv{x} \in
  \R^d$. $o$, $h$: output and hidden dimensions of the deep linear network
  in \cref{sec:dln}.
  \item $\alpha$: the load $P/N$ in Block 1; the spectral-decay exponent
  in \cref{sec:scaling} (both uses are entrenched in their literatures;
  we flag the switch where it happens).
  \item $\varepsilon$: generalization error; $\varsigma^2$: label-noise
  variance.
  \item $\eta$: learning rate; $\tau = 1/\eta$: gradient-flow time
  constant; $B$: batch size.
  \item $\lambda_i$: kernel eigenvalues (with eigenfunctions $\phi_i$,
  target coefficients $v_i$); in \cref{sec:sgd}, $\lambda$ is a Hessian
  eigenvalue (flagged there). $\kappa$: the effective regularization /
  chemical potential of \cref{thm:eigenlearning}.
  \item $\delta$: the ridge parameter in kernel sections; in local
  asymptotic expansions (near $R = 1$ in \cref{sec:ising}) we follow
  custom and reuse $\delta$ for a small deviation --- context
  disambiguates, and the two never meet.
  \item $\gamma$: the output width-scaling exponent of \cref{sec:gamma};
  $\gamma_0$: the DMFT richness parameter of \cref{sec:dmft} --- related
  in spirit, distinct in definition.
  \item $\nu$: the uniform probability measure on the sphere (Block 1);
  $\mu^N$, $\mu_t$: the empirical/limiting measure of neurons in
  \cref{sec:mup}; $\pi$: the input-data probability measure; $\mu$ as an
  \emph{index} labels training examples in Block 1.
  \item $\chi$: the correlation-map slope (\cref{sec:criticality}), also
  controlling gradient growth and the NTK depth recursion.
\end{itemize}
\end{callout}

%===============================================================================
\section{Session 0: Reading and discussion}
\label{sec:reading}

The assigned reading, \S2 of \cite{simon2026scientific}, argues that a
scientific theory of deep learning --- a ``mechanics of learning'' --- is
already emerging, and presents five converging bodies of evidence:
analytically solvable settings, insightful limits, simple empirical laws,
disentangled hyperparameters, and universal phenomena. The rest of these
notes develops the first two in detail and samples the remaining three. The
purpose of the opening session is to put the paper's \emph{framing} on
trial before we invest a day in its program.

\emph{Format} (70 minutes): 40 minutes reading \S2 individually (students
who read in advance re-skim and prepare positions); 20 minutes discussion
in small groups of 3--4, each group working through the three prompts
below; 10 minutes plenary, with one representative per group reporting the
group's most contested claim.

The starting observation deserves to be stated carefully, because it is easy
to nod along without noticing how strong it is. A deep learning system is
completely specified by four explicit ingredients: an architecture (a
composition of simple linear and nonlinear maps), a dataset, an objective,
and a learning rule (gradient-based updates from a random initialization).
Nothing is hidden. Every parameter, every gradient, every intermediate
activation can be measured at every step, at whatever precision we like. In
the language of the reading: deep learning directly exposes its equations of
motion; \emph{the central challenge is not opacity but complexity}. Contrast
genuinely opaque sciences: in neuroscience or cell biology the equations of
motion are unknown and the state is largely unmeasurable, and in
experimental physics we often infer microscopic laws from scattering
amplitudes. Deep learning is closer to fluid dynamics: we possess the
microscopic law (the update equation is five lines of code) and lack the
effective description.

One refinement of the framing is worth planting before the groups start:
the data distribution is arguably \emph{not} explicit --- we have samples,
not the underlying probability measure. The defensible version of the claim
is that the \emph{trained system} is fully measurable while the
\emph{environment} is known only empirically; this foreshadows the closing
discussion of universal structure in data (\cref{sec:closing}).

The three discussion prompts, with notes on where each tends to lead:

\begin{enumerate}
  \item \emph{Rank the physics analogies in Table 1 of the reading}
  (solvable models $\leftrightarrow$ hydrogen atom; limits
  $\leftrightarrow$ thermodynamic limit; empirical laws $\leftrightarrow$
  Kepler and Planck; hyperparameters $\leftrightarrow$ Reynolds number;
  universality $\leftrightarrow$ critical phenomena). Which are
  load-bearing --- carrying real technical machinery across --- and which
  are decorative? Ask each group to produce an explicit ordering and defend
  its top and bottom choices. Defensible answers vary, but the limits and
  universality analogies tend to survive scrutiny best; by the end of the
  day, students can revisit their ranking with evidence.
  \item \emph{Discretization and other approximations.} The reading's
  Discretization Hypothesis proposes that practical networks are best
  understood as noisy, finite approximations to models of infinite size ---
  like a finite-difference discretization of a PDE --- so that finite-size
  effects typically \emph{cost} performance rather than enable it. Two
  tasks for the group. (a) Inventory the approximations the day ahead will
  use: infinite width, infinite depth, infinite data, continuous time
  (gradient flow in place of discrete SGD steps), population loss in place
  of empirical loss. Which of these are ``discretization-like'' (the finite
  object is a noisy stand-in for a cleaner infinite one) and which might
  change the physics qualitatively? (b) The hypothesis is falsifiable: it
  would be refuted by a finite-size effect that delivers a general benefit
  unachievable in any limit. Collect candidate falsifiers (SGD noise as a
  regularizer? capacity constraints preventing memorization?) and park
  them; we revisit after Blocks 2--3, which supply the limits the
  hypothesis refers to.
  \item \emph{What actually matters for alignment?} Two sub-questions.
  (a) The reading wants falsifiable, quantitative, \emph{average-case}
  predictions rather than worst-case bounds --- physics rather than
  classical learning theory. What is lost? Worst-case guarantees are what
  you ultimately want for adversarial robustness and for safety cases; an
  average-case mechanics of learning complements rather than replaces
  them, and groups should articulate where each is the right tool
  (forecasting typical capabilities vs.\ certifying the absence of rare
  behaviors). (b) Rank the five lines of evidence by expected relevance to
  alignment --- to forecasting what a larger model will do, and to
  monitoring what is forming inside one during training --- and record the
  guesses; the closing session (\cref{sec:closing}) revisits them with the
  day's evidence in hand.
\end{enumerate}

%===============================================================================
\section{Block 1: Analytically solvable settings}
\label{sec:block1}

The oldest and most complete success of physics methods in learning theory
is the statistical mechanics of simple learning machines, developed in the
late 1980s and 1990s \cite{gardner1988space,seung1992statistical,
engel2001statistical}. We treat it first not for historical piety but
because it establishes the \emph{template} that every later section
instantiates: identify an ensemble, find the order parameters, compute a
free entropy by an asymptotic (saddle-point) method, and read off
macroscopic behavior --- including phase transitions --- from where that
free entropy is maximized. All of these terms are defined below; no prior
exposure to statistical mechanics is assumed, and readers who have it
should still skim the dictionary in \cref{sec:dictionary}, which fixes how
each physics word translates into a statement about learning.

\subsection{Statistical mechanics of learning: the teacher--student perceptron}
\label{sec:statmech}

\subsubsection{The setting}

In this section $N$ is the input dimension. A \emph{perceptron} with weight
vector $\bv{w} \in \R^N$ computes the binary function
$\sigma_{\bv{w}}(\bv{x}) = \mathrm{sign}(\bv{w}\cdot\bv{x})$. We study the
\emph{teacher--student} setup: a fixed teacher $\bv{w}^{*}$ generates labels
\begin{equation}
  \label{eq:teacher-data}
  y^\mu = \mathrm{sign}(\bv{w}^{*}\cdot\bv{x}^\mu),
  \qquad \bv{x}^\mu \sim \mathcal{N}(0, I_N) \text{ i.i.d.},
  \qquad \mu = 1,\dots,P,
\end{equation}
and a student $\bv{w}$ must infer the rule from the $P$ examples.

\emph{Why constrain weights to a sphere?} The function
$\mathrm{sign}(\bv{w}\cdot\bv{x})$ depends only on the \emph{direction} of
$\bv{w}$: rescaling $\bv{w} \mapsto \lambda\bv{w}$ changes nothing. The
radius is pure gauge, and fixing it removes a flat direction that would
otherwise make volumes ill-defined. We fix $|\bv{w}|^2 = |\bv{w}^*|^2 = N$
(the \emph{spherical} perceptron) rather than $1$ so that individual
components are $O(1)$ --- the natural scale if one imagines each weight
initialized independently, as in every later section --- and so that inner
products like $\bv{w}\cdot\bv{w}^*$ are extensive (proportional to $N$),
making ``per-coordinate'' quantities finite in the limit.

\emph{The thermodynamic limit.} In physics, the thermodynamic limit sends
the number of degrees of freedom to infinity while holding fixed all
\emph{intensive} quantities (densities, ratios --- quantities that do not
grow with system size); it is the limit in which fluctuations of
macroscopic observables vanish and sharp laws emerge. Here the degrees of
freedom are the $N$ weight components, and the relevant intensive quantity
is the \emph{load}
\begin{equation}
  \label{eq:load}
  \alpha = P/N,
\end{equation}
the number of examples \emph{per parameter} --- ``load'' because it measures
how heavily the data constrains the model. So the thermodynamic limit is
$P, N \to \infty$ at fixed $\alpha$. Why this coupled scaling? If $P$ grew
slower than $N$ ($\alpha \to 0$), each parameter would be asymptotically
unconstrained and no learning would survive the limit; faster
($\alpha \to \infty$), the problem would be asymptotically fully determined
and learning would be trivially perfect. Only at fixed $\alpha$ does a
nontrivial competition survive, with $\alpha$ as the knob. (Compare: all of
\cref{sec:scaling} on scaling laws is about joint limits of data and model
size; the habit starts here.)

\emph{The order parameter.} An \emph{order parameter} is a single
macroscopic quantity that summarizes the microscopic state for the purposes
of the questions asked --- the physicist's compression of $N$ numbers into
one. Here it is the \emph{overlap}
\begin{equation}
  \label{eq:overlap}
  R = \frac{\bv{w}\cdot\bv{w}^{*}}{N} \in [-1, 1],
\end{equation}
the cosine of the angle between student and teacher. The precise sense in
which ``only $R$ matters'' is a symmetry argument: the input distribution
$\mathcal{N}(0, I_N)$ and the uniform measure on the sphere are invariant
under the orthogonal group $O(N)$, so any quantity defined by averaging
over inputs --- the generalization error below, the volume computations to
come --- is invariant under simultaneous rotation of $(\bv{w}, \bv{w}^*)$.
The complete list of rotation invariants of the pair is $|\bv{w}|^2$,
$|\bv{w}^*|^2$, and $\bv{w}\cdot\bv{w}^*$; the first two are fixed by the
spherical constraint, leaving $R$. Any rotation-invariant function of
$(\bv{w},\bv{w}^*)$ is a function of $R$ alone.

\emph{Generalization error.} Define the generalization error of a student
as the probability of disagreeing with the teacher on a fresh input:
\begin{equation}
  \label{eq:gen-error-def}
  \varepsilon(\bv{w}) = \Pr_{\bv{x}\sim\mathcal{N}(0,I_N)}
  \bigl[\mathrm{sign}(\bv{w}\cdot\bv{x}) \neq
  \mathrm{sign}(\bv{w}^*\cdot\bv{x})\bigr].
\end{equation}
By the symmetry argument this is a function of $R$ alone, and it has a
clean geometric form. Project the Gaussian input onto the plane spanned by
$\bv{w}$ and $\bv{w}^*$: the projection is an isotropic two-dimensional
Gaussian (every marginal of an isotropic Gaussian is isotropic), and the
components of $\bv{x}$ orthogonal to the plane affect neither sign. The two
decision boundaries are lines through the origin of this plane, meeting at
angle $\vartheta = \arccos R$; the student and teacher disagree exactly
when the projected input falls in the double wedge between the lines, which
by isotropy has probability $\vartheta/\pi$:
\begin{equation}
  \label{eq:eps-of-R}
  \varepsilon(R) = \frac{1}{\pi}\arccos R .
\end{equation}
The answer is independent of $N$ because only the two-dimensional
projection ever entered. So the entire learning problem reduces to: how
does the \emph{typical} overlap $R$ grow with the load $\alpha$?

\subsubsection{Gibbs learning, the version space, and a dictionary}
\label{sec:dictionary}

Consider the most naive learning rule compatible with the data:
\emph{Gibbs learning}, i.e.\ draw a student uniformly at random from the
\emph{version space}
\begin{equation}
  \label{eq:version-space}
  V = \bigl\{\bv{w} : |\bv{w}|^2 = N,\;
  y^\mu\,(\bv{w}\cdot\bv{x}^\mu) > 0 \;\; \forall \mu \bigr\},
\end{equation}
the set of students consistent with every example. The central object is
the version-space volume
\begin{equation}
  \label{eq:Z-def}
  Z = \int d\nu(\bv{w}) \prod_{\mu=1}^{P}
  \Theta\bigl(y^\mu\, \bv{w}\cdot\bv{x}^\mu\bigr),
\end{equation}
where $d\nu$ is the uniform probability measure on the sphere and $\Theta$
the step function ($\Theta(u) = 1$ for $u > 0$, else $0$). Note that the
product of $\Theta$'s is exactly the indicator function of $V$: the
integrand is the (unnormalized) density of the probability measure we care
about --- the uniform measure on consistent students, which is also the
Bayesian posterior over $\bv{w}$ given the data, under a uniform prior and
a noiseless likelihood.

This is the point to fix the dictionary, in terms a learning theorist
cares about:
\begin{itemize}
  \item \emph{Partition function} $Z$: the normalizing constant of the
  posterior. In Bayesian language, the \emph{evidence} (marginal
  likelihood) of the data. In this zero-noise problem it is simply the
  fraction of hypothesis space consistent with the data.
  \item \emph{Energy}: minus the log-likelihood of a hypothesis. Here the
  likelihood is $0$ or $1$ (an example is classified correctly or not),
  so energy counts violated examples, and only zero-energy hypotheses
  survive: we are at \emph{zero temperature}, meaning constraints are hard.
  At positive temperature one would tolerate violations at a price ---
  i.e., a smooth loss.
  \item \emph{Entropy}: the log of the volume of hypotheses with a given
  property (e.g., a given $R$). Large entropy $=$ many hypotheses.
  \item \emph{Gibbs measure / posterior}: the probability measure
  $\propto e^{-\text{energy}}$ on hypotheses; here, uniform on $V$.
\end{itemize}
Define
\begin{equation}
  \label{eq:free-entropy-def}
  F = -\log Z, \qquad
  \Phi = \frac{1}{N}\,\E_{\mathrm{data}} \log Z = -\frac{1}{N}\,
  \E_{\mathrm{data}} F .
\end{equation}
Conventions, since the signs trip everyone once: physicists' \emph{free
energy} is $F$ (they minimize it; small $F$ means large accessible volume,
which is why the negative sign is wanted --- $F$ decreases as
$Z$ grows); because our $Z$ is a pure volume with no energy scale, $\log Z$
is the log of a volume, i.e.\ an entropy, and $\Phi$ is called the
\emph{free entropy} (to be maximized). We normalize by $1/N$ because
$\log Z$ is \emph{extensive} --- proportional to $N$, being a log-volume in
$N$ dimensions --- and dividing by $N$ produces an intensive,
per-parameter quantity with a finite limit. Finally, \emph{quenched}: the
average $\E_{\mathrm{data}}$ is taken over datasets that are drawn once and
then \emph{frozen} while the student adapts to them (physics term for
randomness fixed before the dynamics run, as impurities are frozen into an
alloy when it is quenched in water). The dataset plays the role that random
couplings play in disordered magnets, and the quenched free entropy $\Phi$
is the correct, typical-case object: $\log Z$ concentrates around
$N\Phi$ for large $N$, so $\Phi$ describes what happens on a typical
dataset.

\subsubsection{The annealed calculation, in full}
\label{sec:annealed}

Averaging $\log Z$ is hard (the log of a sum does not factorize); averaging
$Z$ is easy. The \emph{annealed approximation} swaps them:
\begin{equation}
  \label{eq:annealed-swap}
  \E \log Z \;\longrightarrow\; \log \E Z .
\end{equation}
Let us do the easy computation completely, because its structure recurs all
day. Since the examples are i.i.d., the average of \cref{eq:Z-def} over the
data factorizes across examples; and each example is satisfied by a student
at overlap $R$ with probability $1 - \varepsilon(R)$ (the probability that
teacher and student \emph{agree} on a random input --- the complement of the
wedge event behind \cref{eq:eps-of-R}). Hence
\begin{equation}
  \label{eq:EZ}
  \E Z = \int d\nu(\bv{w})\; \bigl(1 - \varepsilon(R(\bv{w}))\bigr)^{P}.
\end{equation}
Next, trade the $N$-dimensional integral over $\bv{w}$ for a
one-dimensional integral over the collective coordinate $R$: this works
because the integrand depends on $\bv{w}$ only through $R$, so we need only
the volume of each level set. Decompose $\bv{w} = R\,\bv{w}^* +
\bv{w}_\perp$ with $\bv{w}_\perp \perp \bv{w}^*$; the spherical constraint
gives $|\bv{w}_\perp|^2 = N(1 - R^2)$, so the slice at fixed $R$ is a
sphere of radius $\sqrt{N(1-R^2)}$ in the $(N-1)$-dimensional orthogonal
complement, with surface volume $\propto \bigl(N(1-R^2)\bigr)^{(N-2)/2}$.
Therefore
\begin{equation}
  \label{eq:coarea}
  d\nu(\bv{w}) \;=\; e^{N s(R)\, +\, O(\log N)}\,dR,
  \qquad
  s(R) = \tfrac12 \log\bigl(1 - R^2\bigr) + \text{const}.
\end{equation}
``To exponential accuracy'' means precisely this: we keep the term of order
$N$ in the exponent and drop terms of order $\log N$ and below (Jacobians,
prefactors, the surface-area constant); the small parameter of the
expansion is $1/N$, and the dropped terms change $\frac1N \log(\cdot)$ by
$O(\log N / N) \to 0$. Substituting,
\begin{equation}
  \label{eq:annealed-integral}
  \E Z \;\asymp\; \int_{-1}^{1} dR\;
  e^{\,N\,G_{\mathrm{ann}}(R)},
  \qquad
  G_{\mathrm{ann}}(R) = \tfrac12\log(1-R^2)
  + \alpha \log\bigl(1 - \varepsilon(R)\bigr).
\end{equation}
Now \emph{Laplace's method}: for an integral $\int e^{N G(R)}\,dR$ with $N
\to \infty$, the integrand is exponentially peaked at the maximizer $R_*$
of $G$; the region within $\sim 1/\sqrt{N}$ of $R_*$ contributes everything
(expanding $G$ to second order gives a Gaussian of width $1/\sqrt{N\,
|G''|}$), and
\begin{equation}
  \label{eq:laplace}
  \frac{1}{N}\log \int e^{N G(R)}\, dR
  \;\xrightarrow{N\to\infty}\; \max_R G(R).
\end{equation}
The large parameter is $N$; this is the same logic as ``saddle-point
evaluation'' elsewhere in these notes (for complex integrands the maximum
becomes a saddle point in the complex plane, whence the name). Applying it:
\begin{equation}
  \label{eq:annealed-saddle}
  \frac{1}{N}\log \E Z
  = \max_{R}\; G_{\mathrm{ann}}(R) + o(1).
\end{equation}

The two terms of $G_{\mathrm{ann}}$ stage an \emph{entropy--energy
competition}, the basic plot of statistical mechanics. Intuitively: the
first term counts \emph{how many} students sit at overlap $R$ --- it is
maximized at $R = 0$, because in high dimension almost all of the sphere is
nearly orthogonal to any fixed direction (concentration of measure) --- and
it acts as a pull toward the typical, the numerous, the disordered. The
second term measures \emph{how well} such students fit the data --- it is
maximized at $R = 1$ and acts as a pull toward the special, the accurate,
the ordered. Neither wins outright; the compromise moves with $\alpha$,
which weights the data term. In physics the same structure appears with
energy in place of data-fit and temperature in place of $1/\alpha$; ordered
phases (magnets, crystals) form when the energy term wins. Here: for small
$\alpha$ the maximum sits near $R = 0$; as $\alpha$ grows it moves ---
smoothly, for this model --- toward $1$. \Cref{fig:gann} shows the
landscape $G_{\mathrm{ann}}(R)$ for a series of loads $\alpha \in [0, 1]$:
a single maximum, drifting continuously --- keep this picture in mind as
the contrast case when the Ising landscape of \cref{sec:ising} develops a
second, competing basin.

\begin{figure}[ht]
  \centering
  \includegraphics[width=0.62\linewidth]{fig/spherical-Gann.pdf}
  \caption{The annealed free-entropy landscape of the spherical
  teacher--student perceptron, \cref{eq:annealed-integral}, for
  $\alpha \in \{0, 0.25, 0.5, 0.75, 1\}$ (top to bottom), with the
  maximizer $R^*(\alpha)$ marked on each curve. The entropy term pins the
  maximum at $R = 0$ when $\alpha = 0$; increasing the load slides it
  smoothly to larger overlap. There is a single basin at every $\alpha$
  and the maximizer moves continuously: no phase transition. Contrast the
  Ising perceptron of \cref{sec:ising}, where a second basin at $R = 1$
  competes with the interior one and dominance changes hands abruptly.}
  \label{fig:gann}
\end{figure}

The large-$\alpha$ asymptotics can be read off by expanding near $R = 1$.
Writing $R = \cos(\pi\varepsilon)$ so that $\varepsilon$ is the
generalization error itself, $1 - R^2 \approx (\pi\varepsilon)^2$, and
\begin{equation}
  \label{eq:Gann-large-alpha}
  G_{\mathrm{ann}} \approx \log(\pi\varepsilon) - \alpha\,\varepsilon
  + \text{const},
\end{equation}
whose maximum is at
\begin{equation}
  \label{eq:annealed-eps}
  \varepsilon_{\mathrm{ann}}(\alpha) \simeq \frac{1}{\alpha}.
\end{equation}
Learning curves decay as a \emph{power law} in the amount of data, with
exponent $1$ --- our first derived, not fitted, scaling law.

\subsubsection{Interlude: phase transitions, and where they hide here}
\label{sec:phase-transitions}

A \emph{phase transition} is a point in a control parameter (here
$\alpha$; in physics, temperature or pressure) at which the limiting free
entropy $\Phi(\alpha)$ fails to be analytic, so that some macroscopic
observable changes abruptly. The standard taxonomy, translated to our
setting:
\begin{itemize}
  \item \emph{First-order transition}: $\Phi$ is continuous but its first
  derivative $\partial\Phi/\partial\alpha$ jumps. Mechanism: two competing
  local maxima (``basins'') of the free-entropy landscape --- two
  qualitatively different families of students, say poorly-aligned and
  well-aligned --- whose free entropies $\Phi_1(\alpha)$, $\Phi_2(\alpha)$
  cross at $\alpha_c$. The physical branch is the larger one, so
  $\Phi = \max(\Phi_1, \Phi_2)$ has a kink, and the order parameter $R$
  (which differs between the basins) jumps discontinuously.
  \item \emph{Second-order (continuous) transition}: $\partial\Phi/
  \partial\alpha$ is continuous but the second derivative diverges; the
  order parameter is continuous but its derivative is not. No competing
  basins --- one basin deforms singularly.
  \item \emph{Finite-size smoothing}: at finite $N$, $\Phi_N(\alpha)$ is
  analytic --- there are no true singularities in finite systems. A
  first-order jump in $\partial\Phi/\partial\alpha$ appears as a steep
  sigmoid of width $\sim 1/N$, and correspondingly $\partial^2\Phi/
  \partial\alpha^2$ shows a peak whose height grows with $N$ and diverges
  in the limit. Watching such a peak sharpen as $N$ grows is the standard
  numerical diagnostic of a transition --- and the smoothed-sigmoid picture
  is worth remembering when the phase-transitions day discusses
  ``emergence'' in models of growing size.
\end{itemize}

\emph{Is there a transition in our model?} For the spherical teacher--student
perceptron with Gibbs learning: \emph{no}. The function
$G_{\mathrm{ann}}(R)$ (and its quenched counterpart below) has a single
maximum that moves smoothly with $\alpha$; $\Phi(\alpha)$ is analytic, the
learning curve $\varepsilon(\alpha)$ is smooth, and learning is gradual.
The reason is visible in \cref{eq:annealed-integral}: the entropy
$\tfrac12\log(1-R^2)$ diverges to $-\infty$ as $R \to 1$, so the perfectly
aligned state is infinitely penalized and the maximum always sits in the
interior, sliding continuously. To produce a transition we need a model in
which the perfect state carries \emph{finite} entropy --- which is exactly
what discreteness provides.

\subsubsection{A first-order transition, fully worked: the Ising perceptron}
\label{sec:ising}

Change one ingredient: let both teacher and student have \emph{binary}
weights,
\begin{equation}
  \label{eq:ising-setup}
  \bv{w}, \bv{w}^* \in \{\pm 1\}^N
\end{equation}
(the name is by analogy with Ising spins, the $\pm 1$ magnetic moments of
the simplest lattice magnet; no other physics of that model is imported).
The data are as in \cref{eq:teacher-data}, the sign rule and the
generalization error are unchanged --- for any fixed $\bv{w}$ with
$|\bv{w}|^2 = N$, binary or not, the pair $(\bv{w}\cdot\bv{x},
\bv{w}^*\cdot\bv{x})/\sqrt{N}$ is jointly Gaussian with correlation $R$,
so the wedge formula \cref{eq:eps-of-R} still gives
$\varepsilon(R) = \arccos(R)/\pi$. Only one thing changes: the
\emph{entropy of the slice at overlap $R$} is now a counting problem
rather than a surface-area problem. A student agrees with the teacher on a
fraction $p = (1+R)/2$ of its $N$ coordinates, so the number of students
at overlap $R$ is $\binom{N}{pN}$, and Stirling's formula gives
\begin{equation}
  \label{eq:ising-entropy}
  s_{\mathrm{I}}(R) = \frac1N \log\binom{N}{\frac{(1+R)}{2}N}
  \;\longrightarrow\;
  H\!\left(\frac{1+R}{2}\right),
  \qquad
  H(p) = -p\log p - (1-p)\log(1-p),
\end{equation}
the binary entropy function. Two features to stare at, because they cause
everything: $s_{\mathrm{I}}(0) = \log 2$ (the bulk of the hypercube), and
$s_{\mathrm{I}}(1) = 0$ --- \emph{finite}, not $-\infty$: the perfectly
aligned state is a single, legitimate configuration (the teacher itself),
not an infinitely thin shell.

The annealed free entropy is assembled exactly as in
\cref{eq:annealed-integral}:
\begin{equation}
  \label{eq:ising-G}
  G_{\mathrm{I}}(R) = H\!\left(\frac{1+R}{2}\right)
  + \alpha \log\bigl(1 - \varepsilon(R)\bigr),
  \qquad
  \Phi_{\mathrm{ann}}(\alpha) = \max_{R \in [0,1]} G_{\mathrm{I}}(R).
\end{equation}
Now analyze the competition --- this can be done on one blackboard.
\begin{itemize}
  \item \emph{The teacher state.} At $R = 1$: entropy $0$, and
  $\varepsilon(1) = 0$ so the data term is also $0$. Hence
  $G_{\mathrm{I}}(1) = 0$ \emph{for every} $\alpha$: the perfect state is
  always available, at exactly zero free entropy.
  \item \emph{Near $R = 1$}, with $\delta = 1 - R$: the entropy gained by
  deviating is $\sim \tfrac{\delta}{2}\log(1/\delta)$, but the data cost is
  $\alpha\,\varepsilon \sim \alpha\sqrt{2\delta}/\pi$ (expand $\arccos$),
  and $\sqrt{\delta} \gg \delta\log(1/\delta)$: the cost always wins.
  So $R = 1$ is a \emph{local} maximum of $G_{\mathrm{I}}$ at every load
  --- there is a basin around the teacher, separated from the rest.
  \item \emph{The interior branch.} At small $\alpha$ the entropy dominates
  and $G_{\mathrm{I}}$ has its global maximum at an interior $R^*(\alpha)$
  with $G_{\mathrm{I}}(R^*) > 0$ (at $\alpha = 0$ it is $\log 2$ at
  $R = 0$). As $\alpha$ grows, this maximum drifts up in $R$ and
  \emph{down} in value.
  \item \emph{The crossing.} The transition occurs when the two basins
  exchange dominance: $G_{\mathrm{I}}(R^*(\alpha_c)) = 0 =
  G_{\mathrm{I}}(1)$. Solving numerically (a ten-line exercise, see
  \cref{ex:ising}): $\alpha_c^{\mathrm{ann}} \approx 1.45$, with the
  interior maximum still sitting at $R^* \approx 0.70$, i.e.\
  generalization error $\varepsilon^* \approx 0.25$. At that load the
  typical student jumps from $25\%$ error to \emph{zero} error
  discontinuously: a first-order transition to perfect generalization.
  \item \emph{The thermodynamic signature.} On the interior branch,
  $\partial\Phi_{\mathrm{ann}}/\partial\alpha =
  \log(1 - \varepsilon(R^*)) \approx -0.29$; on the teacher branch it is
  $0$. So $\Phi_{\mathrm{ann}}(\alpha) = \max(\,G_{\mathrm{I}}(R^*(\alpha)),
  0\,)$ has a kink at $\alpha_c$: the first derivative jumps by
  $\approx 0.29$, and the second derivative acquires a delta-function
  spike --- at finite $N$, a sigmoidal crossover in
  $\partial\Phi/\partial\alpha$ of width $\sim 1/N$ and a peak in
  $\partial^2\Phi/\partial\alpha^2$ whose height grows with $N$, exactly
  the taxonomy above.
  \item \emph{Metastability.} Past $\alpha_c$ the interior maximum survives
  as a \emph{local} maximum --- in the annealed theory, out to a spinodal
  at $\alpha \approx 1.73$, where it disappears entirely. In the
  metastable window, learning dynamics that improve alignment locally
  (Gibbs sampling, greedy flips) remain stuck at $\varepsilon \approx
  0.25$ even though the perfect state is thermodynamically dominant: the
  information to identify the teacher is present in the data before the
  system can find it. More data eventually helps \emph{abruptly and later
  than it ``should''} --- the prototype for every later conversation about
  capabilities that arrive suddenly and lag their information-theoretic
  possibility.
\end{itemize}
The quenched (replica) treatment refines the numbers but not the picture:
the true transition is at $\alpha_c \approx 1.245$ \cite{gyorgyi1990first}
--- annealing overestimates the transition point for the same reason it
overestimated $\varepsilon$ in \cref{eq:annealed-eps} --- while the
first-order character, the kink, and the metastable window all survive.
Note also the instructive pathology: for $\alpha$ large the annealed
$G_{\mathrm{I}}(R^*)$ would go \emph{negative} on the interior branch,
i.e.\ predict ``fewer than one'' imperfect student --- the negative-entropy
nonsense promised earlier, resolved precisely by recognizing that the
$R = 1$ state has taken over.

\begin{exercise}[The Ising perceptron transition]
\label{ex:ising}
\important
(a) Derive \cref{eq:ising-entropy} from Stirling's formula.
(b) Show that $R = 1$ is a local maximum of $G_{\mathrm{I}}$ for every
$\alpha$, by the $\sqrt{\delta}$ vs.\ $\delta\log(1/\delta)$ comparison.
(c) Numerically compute $\Phi_{\mathrm{ann}}(\alpha)$, locate
$\alpha_c^{\mathrm{ann}}$, plot the order parameter $R^*(\alpha)$ (with its
jump), $\partial\Phi/\partial\alpha$ (with its kink), and find the spinodal
where the interior maximum disappears.
(d) Simulate: for $N \in \{50, 100, 200\}$, run Gibbs sampling (Metropolis
flips of single weights, accepting moves that keep all examples satisfied)
and plot the sampled $\varepsilon$ against $\alpha$. Watch the transition
smooth into a sigmoid at small $N$, sharpen as $N$ grows, and \emph{lag}
the theoretical $\alpha_c$ --- the metastability of part (c) in action.
\end{exercise}

\emph{The storage problem}, briefly, supplies the other kind of
non-analyticity: with random labels and no teacher, the question is not
generalization but whether \emph{any} consistent student exists. The
version-space volume shrinks with $\alpha$ and vanishes at the
\emph{Gardner capacity} --- $\alpha_c = 2$ for the spherical perceptron
\cite{gardner1988space} --- a satisfiability transition ($Z > 0$ to
$Z = 0$).

\begin{callout}[note]
Where are the \emph{deep}-learning phase transitions in this day? Two are
fully worked out in these notes, in genuinely deep or deep-learning-native
settings: the \emph{order/chaos transition} of deep networks at
initialization (\cref{sec:criticality}) --- a continuous transition with
control parameter $\sigma_w^2$, a diverging depth scale, and algebraic
behavior at criticality, whose entire phase diagram is computable --- and
the \emph{interpolation transition} behind double descent
(\cref{sec:eigenlearning} and the lab), where the overfitting factor $E_0$
diverges as $n$ crosses the model's effective capacity. The transitions
most relevant to alignment --- emergence of capabilities during training,
grokking --- have the same mathematical anatomy (competing branches,
order parameters, finite-size smoothing, metastable lag) and get their own
day; the Ising perceptron is deliberately the simplest system in which
that anatomy can be dissected completely.
\end{callout}

\subsubsection{What the annealed approximation gets wrong}

We now return from the phase-transition detour to the spherical model, and
to the annealed learning curve computed in \cref{sec:annealed}. The exact
treatment of that problem
\cite{seung1992statistical,engel2001statistical} yields
\begin{equation}
  \label{eq:quenched-eps}
  \varepsilon_{\mathrm{Gibbs}}(\alpha) \simeq \frac{0.625}{\alpha},
\end{equation}
same exponent as \cref{eq:annealed-eps}, different constant. Why does
annealing err, and in which direction? Start from Jensen's inequality:
$\log \E Z \geq \E \log Z$, always. So the annealed free entropy
\emph{overestimates} the typical log-volume of the version space. To see
what drives the gap, note that $Z$ is a random variable --- it depends on
the dataset --- of the form $Z = e^{N\phi}$ where $\phi$ fluctuates from
dataset to dataset. Exponentiating turns modest fluctuations of $\phi$ into
enormous fluctuations of $Z$: a dataset whose $\phi$ exceeds the typical
value by a constant is exponentially rare, but contributes an exponentially
large $Z$, and the two exponentials compete \emph{at the same order}. The
average $\E Z$ can therefore be dominated by vanishingly rare datasets ---
in our problem, weakly informative datasets that happen to leave an
unusually large version space --- while the typical dataset behaves quite
differently. (This is the same phenomenon as for lognormal random
variables, where the mean can sit far in the upper tail of the
distribution; the mean is a bad summary of a heavy-tailed quantity.)

This also predicts the \emph{direction} of the error, qualitatively.
The rare datasets that dominate $\E Z$ are those with anomalously large
version spaces, i.e.\ datasets that constrain the student unusually weakly;
weakly constrained students sit at lower overlap. The annealed average thus
overweights poorly-aligned students and \emph{overestimates the
generalization error} --- consistently, $1/\alpha > 0.625/\alpha$.

Will the scaling exponent always survive annealing, as it did here? No.
Here the model is realizable (a teacher exists), fluctuations of $\phi$ are
tame, and annealing errs only in the constant. In general the annealed
computation can fail qualitatively --- most famously for models with
discrete weights at large load, where it can predict negative entropies
(fewer than one consistent student, nonsense for a count) and spurious
transition points. The safe summary: annealing is a cheap upper bound on
the free entropy, trustworthy when the version-space volume concentrates,
and diagnostic of its own failure when it produces absurdities. Knowing
when the cheap calculation is enough is part of the craft.

\subsubsection{The replica method, spelled out}
\label{sec:replica-method}

The correct calculation computes $\E \log Z$ directly, via the identity
\begin{equation}
  \label{eq:replica-identity}
  \E \log Z = \lim_{n\to0} \frac{\E Z^n - 1}{n}
  \qquad\text{(from } Z^n = e^{n\log Z} = 1 + n\log Z + O(n^2)\text{)}.
\end{equation}
The strategy: compute $\E Z^n$ for \emph{integer} $n$, where it has a
transparent meaning, obtain a formula analytic in $n$, and continue it to
$n \to 0$. For integer $n$,
\begin{equation}
  \label{eq:EZn}
  \E Z^n = \E \int \prod_{a=1}^{n} d\nu(\bv{w}^a)
  \prod_{\mu=1}^{P} \prod_{a=1}^{n}
  \Theta\bigl(y^\mu\, \bv{w}^a \cdot \bv{x}^\mu\bigr):
\end{equation}
$n$ independent copies (\emph{replicas}) of the student, all judged against
the \emph{same} dataset. Averaging over the shared data is what couples
them. Each example $\mu$ couples to the replicas only through the $n + 1$
scalars $u^a = \bv{w}^a\cdot\bv{x}^\mu/\sqrt N$ and $v =
\bv{w}^*\cdot\bv{x}^\mu/\sqrt N$, which are jointly Gaussian with
covariance determined entirely by the \emph{overlap matrix}
\begin{equation}
  \label{eq:overlap-matrix}
  q^{ab} = \frac{\bv{w}^a \cdot \bv{w}^b}{N},
  \qquad
  R^a = \frac{\bv{w}^a \cdot \bv{w}^*}{N}
\end{equation}
(the same symmetry argument as before, now for $n+1$ vectors). The data
average factorizes across the $P$ examples, each contributing an identical
factor $e^{\alpha N \mathcal{A}(q, R) / P \cdot P}$; the $\bv{w}$-integrals
at fixed overlaps contribute an entropic factor $e^{N S(q, R)}$ by the same
slice-counting as \cref{eq:coarea}. Laplace's method then gives a
saddle-point problem over the matrix $(q^{ab}, R^a)$:
\begin{equation}
  \label{eq:replica-saddle}
  \frac1N \log \E Z^n = \mathrm{extr}_{q, R}
  \bigl[\, S(q, R) + \alpha\, \mathcal{A}(q, R) \,\bigr] + o(1).
\end{equation}
The disorder is gone; the price is an interaction \emph{among replicas},
carried by $q^{ab}$.

\emph{Why is $q^{ab} = q$ justified here?} The saddle-point equations are
symmetric under permuting replicas, which makes the \emph{replica-symmetric
(RS) ansatz} $q^{ab} = q$ (all $a \neq b$), $R^a = R$ natural --- but
symmetric equations can have symmetry-broken solutions, so naturalness is
not an argument. The honest justification interprets the order parameter:
at the saddle, $q$ is the typical overlap between \emph{two independent
samples from the posterior} (two students drawn from the version space).
If the version space is one connected, roughly convex blob, two independent
draws land at a characteristic distance and a single value of $q$ describes
them: replica symmetry holds. And for the spherical perceptron the version
space \emph{is} convex --- it is the intersection of half-spaces
($y^\mu\,\bv{w}\cdot\bv{x}^\mu > 0$) with the sphere, and convexity (of the
cone; the spherical slice is geodesically well-behaved) guarantees a single
cluster. This is a theorem-grade reason, and pure mathematicians in the
audience should enjoy that the physicists' ansatz is here underwritten by
convex geometry.

\emph{When does replica symmetry break?} When the posterior fragments into
many well-separated clusters: two draws then have (at least) two
characteristic overlaps --- same-cluster and different-cluster --- and
$q^{ab}$ must take multiple values, in a hierarchical pattern in the worst
case. The physical picture is a \emph{glassy phase}: an energy landscape
with exponentially many near-degenerate local basins separated by high
barriers, as in a \emph{spin glass} --- a magnet with random frustrated
couplings (each pair of spins wants a different relative orientation, and
not all preferences can be satisfied at once), whose low-temperature state
is not one ordered configuration but a fractured multitude. Discrete
weights, unrealizable rules, and hard combinatorial problems all push in
this direction; \cite{zdeborova2016statistical} is the modern guide to both
the physics and its rigorous status.

\emph{Finishing the computation.} Within RS, the functions $S$ and
$\mathcal{A}$ reduce to explicit scalar formulas (Gaussian integrals
involving the tail function $H(x) = \int_x^\infty e^{-z^2/2}\,
dz/\sqrt{2\pi}$), the $n \to 0$ limit is taken, and a further symmetry of
Gibbs learning --- a student drawn uniformly from the version space is
statistically indistinguishable from the teacher, forcing $q = R$ at the
saddle --- reduces everything to one scalar equation. Its numerical solution
gives the full learning curve, and its large-$\alpha$ expansion gives
\cref{eq:quenched-eps}. The guided appendix (\cref{sec:replica-appendix})
walks through these steps; the point of this subsection is that every one
of them is an application of ideas already on the table --- symmetry,
slice-counting, Laplace --- plus one genuinely new and slightly illicit
move, the continuation $n \to 0$.

\begin{callout}[tip]
What the machinery buys, in summary: exact learning curves for a model
class, the entropy--energy mechanism behind them, and --- in the Ising
variant --- a first-order phase transition to perfect generalization at
$\alpha_c \approx 1.245$, with metastability explaining why local learning
dynamics lag the information-theoretic transition. These are the prototypes
for every ``sudden capability'' discussion later in the course.
\end{callout}

\begin{exercise}[Annealed learning curve]
\label{ex:annealed}
Fill in the steps of the annealed calculation: derive
\cref{eq:eps-of-R}, the sphere-slice entropy $s(R)$, and the saddle-point
condition of \cref{eq:annealed-saddle}; verify
\cref{eq:annealed-eps}. Then show that at small $\alpha$ the annealed theory
predicts $R \propto \alpha$, and sketch the full annealed learning curve.
\end{exercise}

\subsection{Deep linear networks: how the solvable dynamics arise}
\label{sec:dln}

Students met deep linear networks on an earlier day; this section is
written to be followable if that day has faded, at the cost of some
repetition. The pattern to extract --- \emph{dynamics decouple in the right
basis, each mode is solvable, modes are learned greedily} --- recurs in the
NTK (\cref{sec:ntk}) and quietly underlies the developmental-stages
material of the SLT day.

\subsubsection{Setup and gradient flow}

A depth-2 \emph{linear} network computes $f(\bv{x}) = W_2 W_1 \bv{x}$,
where inputs are $\bv{x} \in \R^d$, targets $\bv{y} \in \R^{o}$, and the
weight matrices are $W_1 \in \R^{h \times d}$, $W_2 \in \R^{o \times h}$
($d$: input dimension; $h$: hidden width; $o$: output dimension). There is
no nonlinearity, so
the network can only implement a linear map --- but the \emph{training
dynamics} of $(W_1, W_2)$ are genuinely nonlinear, and that is the point.
Train by \emph{gradient flow} (the continuous-time limit of gradient
descent, $\dot\theta = -\eta\,\nabla_\theta\mathcal{L}$; we write
$\tau = 1/\eta$ for the resulting time constant, so that large $\tau$
means slow learning) on the population quadratic loss
\begin{equation}
  \label{eq:dln-loss}
  \mathcal{L}(W_1, W_2)
  = \tfrac12\, \E\,\bigl|\bv{y} - W_2 W_1 \bv{x}\bigr|^2 ,
\end{equation}
with \emph{whitened} inputs, $\E[\bv{x}\bv{x}^{\top}] = I_d$ (whitening ---
rotating and rescaling inputs so their covariance is the identity ---
removes an inessential complication; it can always be arranged by
preprocessing). Expanding \cref{eq:dln-loss} and using whitening, the data
enter the gradients only through the \emph{input--output correlation
matrix}
\begin{equation}
  \label{eq:sigma-yx}
  \Sigma^{yx} = \E\bigl[\bv{y}\,\bv{x}^{\top}\bigr] \in \R^{o\times d},
\end{equation}
and gradient flow reads
\begin{equation}
  \label{eq:dln-flow}
  \tau\,\dot W_1 = W_2^{\top}\bigl(\Sigma^{yx} - W_2 W_1\bigr),
  \qquad
  \tau\,\dot W_2 = \bigl(\Sigma^{yx} - W_2 W_1\bigr) W_1^{\top} .
\end{equation}
(Derive these: a two-line matrix-calculus exercise.) Note the structure:
each layer's motion is driven by the common residual $\Sigma^{yx} - W_2
W_1$, filtered through the \emph{other} layer --- the layers are coupled,
which is where all the nonlinearity lives.

\emph{Relation to a teacher matrix.} On the earlier day we assumed
$\bv{y} = M\bv{x}$ for a teacher matrix $M$. With whitened inputs that
gives $\Sigma^{yx} = \E[M\bv{x}\bv{x}^\top] = M$: the input--output
correlation \emph{is} the teacher. The present formulation is strictly more
general --- $\bv{y}$ may be noisy or not linearly generated at all --- and
then $\Sigma^{yx}$ is the best linear map in the least-squares sense
(for whitened inputs), i.e.\ the effective teacher the network can hope to
learn. Everything below depends on the data only through $\Sigma^{yx}$.

\subsubsection{Decoupling in the singular basis}

Take the singular value decomposition
\begin{equation}
  \label{eq:svd}
  \Sigma^{yx} = U S V^{\top},
  \qquad S = \mathrm{diag}(s_1 \ge s_2 \ge \dots \ge 0),
\end{equation}
with $U, V$ orthogonal: $s_k$ measures the strength of the $k$-th
input--output mode of the data (for the teacher case, the singular values
of $M$). Change variables to align the network with this basis:
\begin{equation}
  \label{eq:rotated-vars}
  \bar A = U^{\top} W_2, \qquad \bar B = W_1 V
  \qquad (\text{so } W_2 W_1 = U\,\bar A\,\bar B\,V^{\top}).
\end{equation}
Substituting into \cref{eq:dln-flow} and using orthogonality,
\begin{equation}
  \label{eq:rotated-flow}
  \tau\,\dot{\bar B} = \bar A^{\top}\bigl(S - \bar A \bar B\bigr),
  \qquad
  \tau\,\dot{\bar A} = \bigl(S - \bar A \bar B\bigr)\bar B^{\top} :
\end{equation}
the same equations with $\Sigma^{yx}$ replaced by the diagonal $S$. A
\emph{decoupled (spectral) initialization} is one in which $\bar A$ and
$\bar B$ start diagonal --- equivalently, the initial weights are aligned
with the task's singular vectors, $W_2(0) = U A_0$, $W_1(0) = B_0
V^{\top}$ with $A_0, B_0$ diagonal. Then the right-hand sides of
\cref{eq:rotated-flow} are diagonal, diagonality is preserved for all time,
and the matrix ODEs collapse into independent scalar pairs, one per mode:
writing $a_k, b_k$ for the $k$-th diagonal entries,
\begin{equation}
  \label{eq:mode-odes}
  \tau\,\dot a_k = b_k\,(s_k - a_k b_k),
  \qquad
  \tau\,\dot b_k = a_k\,(s_k - a_k b_k).
\end{equation}
Here $a_k b_k$ is the strength with which the network currently implements
mode $k$, and $s_k$ the strength the data demands. (The spectral
initialization is an analytical idealization: a small \emph{random}
initialization is approximately decoupled, and \cite{saxe2014exact} verify
numerically that its dynamics track the decoupled solution closely. Nothing
below depends delicately on it.)

\subsubsection{Solving one mode}

Subtracting the two equations in \cref{eq:mode-odes} shows
$\tfrac{d}{dt}(a_k^2 - b_k^2) = 0$: the imbalance $a_k^2 - b_k^2$ is a
conserved quantity of the flow (a first sign of the conservation structure
of gradient flows on layered models). On the \emph{balanced} manifold
$a_k = b_k$, the mode strength $u_k = a_k b_k = a_k^2$ closes on itself:
\begin{equation}
  \label{eq:saxe-ode}
  \tau\,\dot u_k = 2\,u_k\,(s_k - u_k),
\end{equation}
a logistic (Bernoulli) equation: growth rate proportional to $u_k$ itself
(small modes grow slowly \emph{because} they are small --- the coupling of
the two layers multiplies the signal by the current mode strength), with
saturation at $u_k = s_k$. Separating variables gives
\begin{equation}
  \label{eq:saxe-solution}
  u_k(t) = \frac{s_k}
  {1 + \bigl(s_k/u_k(0) - 1\bigr)\, e^{-2 s_k t/\tau}} ,
\end{equation}
a sigmoid that lingers near $0$, rises sharply, and saturates at $s_k$
(\cref{fig:modes}). The rise time --- say, time to reach $s_k/2$ --- is
\begin{equation}
  \label{eq:rise-time}
  t_k \simeq \frac{\tau}{2 s_k}\,\log\!\frac{s_k}{u_k(0)} ,
\end{equation}
so \emph{modes are learned in order of singular value}, and the sharpness
of the stagewise structure is controlled by the initialization scale: as
$u_k(0) \to 0$ the trajectory becomes a staircase --- long plateaus during
which the network idles near a saddle point of the loss (a stationary
point that is a minimum in some directions and a maximum in others),
punctuated by rapid drops as each successive mode switches on.

\begin{figure}[ht]
  \centering
  \includegraphics[width=0.62\linewidth]{fig/mode-trajectories.pdf}
  \caption{Sequential learning of singular modes in a depth-2 linear
  network: the exact solution \cref{eq:saxe-solution} for
  $s_k \in \{3, 2, 1\}$, $u_k(0) = 10^{-6}$, $\tau = 1$ (solid), with the
  asymptotes $u = s_k$ (dotted). Larger singular values are learned first,
  and each mode turns on as a sharp sigmoid after a long plateau ---
  saddle-to-saddle dynamics. This reproduces the phenomenon shown in
  Fig.~1a of \cite{simon2026scientific} (there in the form of Fig.~3 of
  \cite{saxe2014exact}); the figure here is generated directly from
  \cref{eq:saxe-solution}.}
  \label{fig:modes}
\end{figure}

\begin{remark}[Random teachers, and the distribution of rise times]
\label{rem:random-teacher}
Suppose the teacher $M$ is itself random --- i.i.d.\ entries of variance
$\varsigma^2/d$, say, with $o = d$ large. Then classical random matrix
theory gives the limiting distribution of its singular values: the
\emph{quarter-circle} density $p(s) = \sqrt{4\varsigma^2 - s^2}\,/\,
(\pi\varsigma^2)$ on $[0, 2\varsigma]$ (the singular-value form of the
Marchenko--Pastur/semicircle laws; rectangular shapes $o/d = c \neq 1$
deform it and can open a gap at $0$). Feed this into
\cref{eq:rise-time}: since $t(s) \propto \log(s/u_0)/s$ diverges as
$s \to 0$ and $p(s)$ stays strictly positive at $s = 0$, the induced
distribution of rise times is heavy-tailed --- up to logarithmic
corrections, $\Pr[t_k > T] \sim p(0)\,\tau\log(\cdot)/(2T) \propto 1/T$,
a Pareto-like tail with density $\sim T^{-2}$ and \emph{no characteristic
timescale}. The aggregate loss, a superposition of sharp per-mode sigmoids
with power-law-distributed onset times, is therefore a smooth, slowly
decaying curve. Hold this thought: a smooth macroscopic learning curve
assembled from many sharp microscopic events, with power-law statistics
inherited from a spectrum, is precisely the structure that reappears in
the scaling-law theory of \cref{sec:scaling}. If instead $\Sigma^{yx}$ is
\emph{estimated} from $P$ samples, it concentrates on $M$ with random
fluctuations of size $\sim\sqrt{d/P}$; modes with $s_k$ below that noise
floor are effectively invisible --- an effective rank cutoff set by the
data budget.
\end{remark}

\subsubsection{Lessons to carry forward}

\begin{itemize}
  \item \emph{Greedy low-rank bias.} The network acquires the largest
  structure in the data first. The effect is strengthened by small
  initialization, greater depth, minibatch noise, and $\ell_2$
  regularization (references collected in \cite{simon2026scientific}).
  \item \emph{Nonlinear in parameters, linear in data.} All of this
  happens with a linear input--output map: the nontrivial dynamics come
  entirely from the parameterization. Feature learning has a
  parameter-space, not function-space, origin.
  \item \emph{Saddle-to-saddle dynamics} is the cleanest provable model of
  ``developmental stages'' in training --- the empirical phenomenon that
  motivates the SLT day.
\end{itemize}

\begin{exercise}[Deep linear warm-up]
\label{ex:dln}
(a) Derive \cref{eq:dln-flow} from \cref{eq:dln-loss}. (b) Verify that
diagonality of $(\bar A, \bar B)$ is preserved by \cref{eq:rotated-flow}
and that $a_k^2 - b_k^2$ is conserved. (c) Solve \cref{eq:saxe-ode} by
separation of variables to obtain \cref{eq:saxe-solution}, and derive the
rise time \cref{eq:rise-time}. (d) For depth $L$ (balanced initialization
throughout), show the mode ODE becomes $\tau\,\dot u_k = L\,
u_k^{2 - 2/L}(s_k - u_k)$ and discuss how depth sharpens the plateaus.
\end{exercise}

\subsection{Infinite width as the free theory: NNGP and NTK}
\label{sec:ntk}

We now build the ``free field theory'' of deep learning: the infinite-width
limit, in which networks become Gaussian processes and training becomes
linear. Everything later in the day is a statement about how real networks
depart from this solvable baseline.

\subsubsection{Interlude: what a Gaussian process is, three ways}
\label{sec:gp-defs}

A network with random weights is a \emph{random function}: fix the
architecture, draw the parameters, and you have drawn an
$f: \R^d \to \R$. The right language for random functions, and the one
this section's main theorem is stated in, is that of Gaussian processes.
We give the definition three times, in increasing order of physics
content; the third introduces path-integral notation used informally
throughout the rest of the notes.

\emph{Definition 1 (finite marginals).} A random function $f$ is a
\emph{Gaussian process (GP)} with mean function $m(\bv{x}) =
\langle f(\bv{x})\rangle$ and covariance (kernel) $K(\bv{x},\bv{x}')$ if
for \emph{every} finite collection of inputs $\bv{x}_1, \dots,
\bv{x}_k$, the random vector $\bigl(f(\bv{x}_1), \dots,
f(\bv{x}_k)\bigr)$ is multivariate Gaussian, with means $m(\bv{x}_a)$ and
covariance matrix $K(\bv{x}_a, \bv{x}_b)$. Informally: an
infinite-dimensional Gaussian vector, indexed by input points instead of
coordinates. (Throughout we write $\langle\,\cdot\,\rangle$ for
expectation over the random function.) All our GPs have $m = 0$, and we
assume this below.

\emph{Definition 2 (connected correlators).} For observables
$\mathcal{O}_1, \mathcal{O}_2, \dots$ (any functions of $f$ ---
evaluations $f(\bv{x}_a)$ are the ones we need), define the
\emph{connected correlation functions} (cumulants) recursively:
$\langle \mathcal{O}_1 \rangle^c = \langle \mathcal{O}_1 \rangle$,
\begin{equation}
  \label{eq:connected-2}
  \langle \mathcal{O}_1 \mathcal{O}_2 \rangle^c
  = \langle \mathcal{O}_1 \mathcal{O}_2 \rangle
  - \langle \mathcal{O}_1 \rangle \langle \mathcal{O}_2 \rangle,
\end{equation}
and in general by declaring that every full correlator is the sum, over
all partitions of its observables, of products of connected correlators of
the blocks:
\begin{equation}
  \label{eq:moments-cumulants}
  \langle \mathcal{O}_1 \cdots \mathcal{O}_k \rangle
  = \sum_{\text{partitions } \Pi \text{ of } \{1,\dots,k\}}
  \;\prod_{\text{blocks } B \in \Pi}
  \bigl\langle \textstyle\prod_{i \in B} \mathcal{O}_i
  \bigr\rangle^{c},
\end{equation}
which determines $\langle \cdots \rangle^c$ order by order (solve for the
top block). The connected correlator measures what a joint expectation
knows beyond all factorizations into smaller pieces. Then: $f$ is a
(mean-zero) GP with kernel $K$ if and only if
\begin{equation}
  \label{eq:gp-cumulants}
  \langle f(\bv{x})\, f(\bv{x}') \rangle^c = K(\bv{x},\bv{x}'),
  \qquad
  \langle f(\bv{x}_1) \cdots f(\bv{x}_k) \rangle^c = 0
  \quad \text{for all } k \ge 3 .
\end{equation}
Feeding \cref{eq:gp-cumulants} back into \cref{eq:moments-cumulants}
yields \emph{Wick's theorem} (Isserlis, to probabilists): every moment of
a GP is a sum over pairings of two-point functions, e.g.\
$\langle f_1 f_2 f_3 f_4 \rangle = K_{12}K_{34} + K_{13}K_{24} +
K_{14}K_{23}$ in obvious shorthand. This definition makes ``how
non-Gaussian is it?'' quantitative: the first nonvanishing higher
connected correlator is the answer --- exactly the diagnostic used at
finite width in \cref{sec:eft}.

\emph{Definition 3 (effective action).} Write expectations over the random
function as a \emph{functional (path) integral}: for a probability density
$P[f] \propto e^{-S[f]}$ over functions,
\begin{equation}
  \label{eq:path-integral}
  \langle \mathcal{O} \rangle
  = \frac{1}{Z} \int \mathcal{D}f\; \mathcal{O}[f]\, e^{-S[f]},
  \qquad
  Z = \int \mathcal{D}f\; e^{-S[f]},
\end{equation}
where $S[f]$ is called the \emph{action}. Connected correlators and the
action generate each other: sourcing the integral,
$Z[J] = \int \mathcal{D}f\, e^{-S[f] + \int J f}$, one has
\begin{equation}
  \label{eq:generating}
  \langle f(\bv{x}_1)\cdots f(\bv{x}_k) \rangle^c
  = \left.\frac{\delta^k \log Z[J]}
  {\delta J(\bv{x}_1) \cdots \delta J(\bv{x}_k)}\right|_{J=0},
\end{equation}
so the expansion of $\log Z[J]$ --- equivalently, the polynomial terms
(``couplings'') in the action --- encode the connected correlators order
by order: quadratic terms carry two-point information, quartic terms
generate connected four-point functions, and so on. A GP is then the case
of a purely \emph{quadratic} action,
\begin{equation}
  \label{eq:gaussian-action}
  S[f] = \tfrac12 \int\!\!\int d\bv{x}\, d\bv{x}'\;
  f(\bv{x})\, K^{-1}(\bv{x},\bv{x}')\, f(\bv{x}'),
\end{equation}
with $K^{-1}$ the \emph{inverse kernel}, i.e.\ the inverse of $K$ as an
integral operator, $\int K^{-1}(\bv{x},\bv{x}'')\,K(\bv{x}'',\bv{x}')\,
d\bv{x}'' = \delta(\bv{x}-\bv{x}')$. In field-theory language: a GP is a
\emph{free field}, interactions are precisely non-Gaussianity, and
``nearly Gaussian'' means small higher couplings. Two caveats for the
mathematicians: at our level of rigor the functional integral is
organized notation, not analysis --- restricted to any finite set of
points, \cref{eq:gaussian-action} is literally the multivariate Gaussian
density with matrix $K^{-1}$, and that finite-dimensional statement is
all we ever use; and $K^{-1}$ may only exist on the appropriate
subspace, which never troubles us for the same reason.

The three definitions are equivalent where they overlap, and each earns
its keep today: Definition 1 is what we verify for networks below;
Definition 2 is how deviations from Gaussianity are measured at finite
width (\cref{sec:eft}); Definition 3 is the dictionary that makes
``infinite width $=$ free theory, $1/N$ corrections $=$ interactions'' a
precise slogan rather than an analogy \cite{roberts2022principles,
halverson2021neural}.

\subsubsection{Networks as Gaussian processes: the recursion, derived}
\label{sec:nngp-derivation}

Now the theorem this language was built for. Consider a fully-connected
network acting on inputs $\bv{x} \in \R^d$ ($d$: the input dimension),
with hidden layers of width $N$. Note the change of cast relative to
\cref{sec:statmech}: from here on, $N$ denotes the \emph{width} of a
hidden layer, and the input dimension gets its own letter. The
preactivations are defined recursively by
\begin{equation}
  \label{eq:preactivations}
  z^{(1)}_i(\bv{x}) = \sum_{j=1}^{d} W^{(1)}_{ij} x_j + b^{(1)}_i,
  \qquad
  z^{(\ell+1)}_i(\bv{x}) = \sum_{j=1}^{N} W^{(\ell+1)}_{ij}\,
  \phi\bigl(z^{(\ell)}_j(\bv{x})\bigr) + b^{(\ell+1)}_i,
\end{equation}
with weights initialized i.i.d.\ $W^{(1)}_{ij} \sim \mathcal{N}(0,
\sigma_w^2/d)$, $W^{(\ell+1)}_{ij} \sim \mathcal{N}(0, \sigma_w^2/N)$
(the $1/(\text{fan-in})$ variance keeps preactivations $O(1)$), and
biases $b^{(\ell)}_i \sim \mathcal{N}(0, \sigma_b^2)$.

\emph{Base case.} Fix any finite collection of inputs. Each
$z^{(1)}_i(\bv{x})$ is a finite linear combination of the Gaussian
variables $W^{(1)}_{ij}, b^{(1)}_i$ with \emph{deterministic}
coefficients $x_j$ --- so the first layer is \emph{exactly} Gaussian at
any width, jointly over inputs and units, no limit needed. Its units are
independent across $i$ (disjoint weight rows), with covariance
\begin{equation}
  \label{eq:K1}
  K^{(1)}(\bv{x},\bv{x}')
  = \sigma_w^2\, K^{(0)}(\bv{x},\bv{x}') + \sigma_b^2,
  \qquad
  K^{(0)}(\bv{x},\bv{x}') := \frac{\bv{x}\cdot\bv{x}'}{d},
\end{equation}
where $K^{(0)}$, the normalized input Gram kernel, is the initial
condition of the whole recursion: the only place the raw data geometry
enters.

\emph{Induction step.} Assume (inductive hypothesis) that as
$N \to \infty$ the layer-$\ell$ preactivations $\{z^{(\ell)}_j\}_j$
approach i.i.d.\ draws, across $j$, of a mean-zero GP with kernel
$K^{(\ell)}$. Condition on layer $\ell$. Then $z^{(\ell+1)}_i(\bv{x})$
in \cref{eq:preactivations} is a linear combination of the fresh
Gaussians $W^{(\ell+1)}_{ij}, b^{(\ell+1)}_i$ with (conditionally)
deterministic coefficients $\phi(z_j^{(\ell)}(\bv{x}))$ --- hence
\emph{conditionally} Gaussian, exactly, with conditional covariance
\begin{equation}
  \label{eq:conditional-cov}
  \mathrm{Cov}\bigl(z^{(\ell+1)}_i(\bv{x}),\,
  z^{(\ell+1)}_i(\bv{x}')\,\big|\, z^{(\ell)}\bigr)
  = \sigma_w^2\, \frac{1}{N} \sum_{j=1}^{N}
  \phi\bigl(z^{(\ell)}_j(\bv{x})\bigr)\,
  \phi\bigl(z^{(\ell)}_j(\bv{x}')\bigr) + \sigma_b^2 ,
\end{equation}
independent across $i$. The random object on the right is an empirical
average of i.i.d.\ quantities (by the inductive hypothesis), so the law
of large numbers replaces it by its expectation as $N \to \infty$:
\begin{equation}
  \label{eq:nngp-recursion}
  K^{(\ell+1)}(\bv{x},\bv{x}') = \sigma_w^2\,
  \E_{(u,v)}\bigl[\phi(u)\phi(v)\bigr] + \sigma_b^2,
  \qquad
  (u,v) \sim \mathcal{N}\!\left(0,
  \begin{pmatrix} K^{(\ell)}(\bv{x},\bv{x}) & K^{(\ell)}(\bv{x},\bv{x}')\\
  K^{(\ell)}(\bv{x},\bv{x}') & K^{(\ell)}(\bv{x}',\bv{x}') \end{pmatrix}
  \right).
\end{equation}
This is the step where randomness dies: the conditional covariance
\emph{self-averages} --- it no longer depends on the realization of layer
$\ell$ --- so conditional Gaussianity becomes unconditional Gaussianity,
with deterministic kernel. Independence across units $i$ survives
(disjoint weight rows again), which is exactly the inductive hypothesis
at layer $\ell + 1$, and the induction closes. (What is swept under the
rug: exchanging the limit with the conditioning requires the
LLN fluctuations, of size $1/\sqrt{N}$, to be controlled uniformly ---
this is where the real proofs \cite{neal1996bayesian,lee2018deep} do
their work, and where the $1/\sqrt N$ will return with a vengeance in
\cref{sec:eft}.)

So the network at initialization is a draw from a Gaussian process with
kernel computed by iterating \cref{eq:nngp-recursion} from the initial
condition \cref{eq:K1} --- the \emph{NNGP kernel}. In the language of
\cref{sec:gp-defs}: the infinite-width network has a quadratic effective
action with kernel $K^{(L)}$, all connected correlators beyond the second
vanish, and everything one can ask about the initialized network is
computable from a two-point function. Note also the structure of
\cref{eq:nngp-recursion}: a deterministic recursion for a two-point
function across depth --- we will meet it again as a dynamical system in
\cref{sec:criticality}.

\begin{exercise}[Gaussian processes, three ways]
\label{ex:gp}
(a) Verify \cref{eq:K1}, including the bias term. (b) From
\cref{eq:gp-cumulants,eq:moments-cumulants}, derive Wick's theorem for
$\langle f_1 f_2 f_3 f_4 \rangle$, and check it directly for a bivariate
Gaussian. (c) For a single hidden layer at \emph{finite} width $N$ (with
$\phi$ bounded), show that the output's connected four-point function is
$O(1/N)$, and identify the $N$-independent Gaussian it decorates. This is
the seed of \cref{sec:eft}.
\end{exercise}

\subsubsection{Linearized training and the neural tangent kernel}

Now train. Expand the network to first order around its initialization:
\begin{equation}
  \label{eq:linearization}
  f_{\mathrm{lin}}(\bv{x};\theta) = f(\bv{x};\theta_0)
  + \nabla_\theta f(\bv{x};\theta_0)^{\top}(\theta - \theta_0).
\end{equation}
This model is \emph{linear in the parameters} but still highly nonlinear
in the input. Write $\psi(\bv{x}) = \nabla_\theta f(\bv{x};\theta_0)$, a
fixed vector for each input. Then \cref{eq:linearization} reads
$f_{\mathrm{lin}} = \text{const} + (\theta - \theta_0)\cdot\psi(\bv{x})$:
the model is a \emph{linear function of a fixed nonlinear embedding} of
the input. In the terminology of kernel methods, such an embedding is
called a \emph{feature map}: it sends each input to a vector of
``features'' (here, the sensitivities of the output to each parameter),
and the learner is only allowed to take linear combinations of them.
Learning a linear model on top of a fixed feature map is \emph{linear
regression} --- convex, exactly solvable --- and everything about it is
controlled by the Gram matrix of the features, i.e.\ the matrix of their
pairwise inner products. That Gram structure gets a name:

\begin{definition}[neural tangent kernel]
\label{def:ntk}
$\displaystyle
  \ntk(\bv{x},\bv{x}') = \psi(\bv{x})^{\top}\psi(\bv{x}')
  = \nabla_{\theta} f(\bv{x};\theta_{0})^{\top}
  \nabla_{\theta} f(\bv{x}';\theta_{0})$.
\end{definition}

\emph{Function-space dynamics, derived.} Fix a training set of $n$ pairs
$(\bv{x}_\mu, y_\mu)$, $\mu = 1, \dots, n$, and write $X =
(\bv{x}_1, \dots, \bv{x}_n)$ for the collection of training inputs and
$\bv{y} = (y_1, \dots, y_n)$ for the targets. For any kernel $K$ we use
the standard block shorthand: $K(X,X)$ is the $n \times n$ Gram matrix
with entries $K(\bv{x}_\mu, \bv{x}_\nu)$, and $K(\bv{x}, X)$ the
$n$-vector with entries $K(\bv{x}, \bv{x}_\mu)$. Train by gradient flow
on the empirical quadratic loss $L(\theta) = \tfrac12\sum_\mu
(f(\bv{x}_\mu;\theta) - y_\mu)^2$:
$\dot\theta = -\nabla_\theta L = -\sum_\mu \nabla_\theta
f(\bv{x}_\mu;\theta)\,(f(\bv{x}_\mu;\theta) - y_\mu)$. What does the
\emph{function} do? Chain rule:
\begin{equation}
  \label{eq:ntk-flow-derivation}
  \partial_t f_t(\bv{x})
  = \nabla_\theta f(\bv{x};\theta_t)^{\top} \dot\theta_t
  = -\sum_{\mu} \nabla_\theta f(\bv{x};\theta_t)^{\top}
  \nabla_\theta f(\bv{x}_\mu;\theta_t)\,
  \bigl(f_t(\bv{x}_\mu) - y_\mu\bigr),
\end{equation}
i.e.\ for \emph{any} network, exactly,
\begin{equation}
  \label{eq:ntk-flow}
  \partial_t f_t(\bv{x}) = -\sum_{\mu}
  K_t(\bv{x}, \bv{x}_\mu)\,\bigl(f_t(\bv{x}_\mu) - y_\mu\bigr),
\end{equation}
where $K_t$ is the tangent kernel of \cref{def:ntk} evaluated at the
\emph{current} parameters $\theta_t$. No approximation has been made:
gradient flow in parameter space is always kernel dynamics in function
space, with a moving kernel. (This is the same statement as
\cref{eq:output-ode} in the mean-field section --- the two normalizations
differ only in what the kernel does next.)

For the \emph{linearized} model the story ends immediately:
$\nabla_\theta f_{\mathrm{lin}}(\bv{x};\theta) = \psi(\bv{x})$ is
independent of $\theta$ by construction, so $K_t = \ntk$ exactly, for all
$t$ and any width. This is what ``training is linear regression, driven
by the NTK at initialization'' means concretely: the residuals
$\Delta_\mu = f_t(\bv{x}_\mu) - y_\mu$ obey the \emph{linear} ODE
$\dot\Delta = -\ntk(X,X)\,\Delta$, identical to gradient flow on a linear
least-squares problem in the features $\psi$; nothing about the network's
nonlinearity in $\bv{x}$ is lost, but all nonconvexity in $\theta$ is
gone. Solving the linear ODE and propagating to a test point:
\begin{equation}
  \label{eq:ntk-solution}
  f_t(\bv{x}) = \ntk(\bv{x},X)\,\ntk(X,X)^{-1}
  \bigl(I - e^{-\ntk(X,X)\,t}\bigr)\bv{y}
  \;\xrightarrow{t\to\infty}\; \text{kernel regression with } \ntk
\end{equation}
(taking $f_0 = 0$ for simplicity). Diagonalizing $\ntk(X,X)$ shows
\emph{spectral bias in time}: the component of the error along the $i$-th
kernel eigenvector decays at rate $\lambda_i$ --- top modes first, exactly
the greedy structure of \cref{eq:saxe-ode}, now for a
nonlinear-in-inputs model.

\subsubsection{Why the kernel freezes: three derivations}
\label{sec:ntk-freeze}

For the linearized model, constancy of the kernel is true by fiat. The
substantive claim of \cite{jacot2018neural} is about the \emph{actual}
network: in the right scaling regime, $K_t \approx K_0$ throughout
training, so the true nonconvex dynamics shadow the linearized ones. Two
miracles are being claimed: $K_0$ concentrates (across random
initializations) on a deterministic kernel, and $K_t$ stays put (across
training time). Neither is automatic --- both fail in mean-field scaling
(\cref{sec:mup}) --- so it is worth having more than one way to see them.

\emph{Derivation 1: displacement power counting.} Work in the
$1/\sqrt{N}$ (NTK) output scaling with width-independent learning rate,
$n$ training points, depth $L$, and $O(1)$ targets. The network has
$\sim N$ parameters per layer (take $N \gg n, L$), and by
\cref{eq:ntk-solution} the function must move an $O(1)$ amount to fit the
data. Because $\ntk = \sum_a \partial_{\theta_a} f\,
\partial_{\theta_a} f$ is a sum of $\sim N$ comparable terms of size
$O(1/N)$ each (that is what the $1/\sqrt N$ scaling arranges), the
function's $O(1)$ motion is assembled from each parameter moving only
$\|\Delta\theta_a\| = O(1/\sqrt{N})$. Now propagate that displacement
into the kernel: the change of $K$ is controlled by second derivatives,
$\Delta K \sim \nabla^2 f \cdot \Delta\theta \cdot \nabla f$, and in this
parameterization the relevant second-derivative contractions are
themselves $O(1/\sqrt N)$; combining, $\Delta K = O(1/N)$ (with depth
factors promoting this to $O(L/N)$ --- the same effective coupling as
\cref{sec:eft}). The regime, stated in full: \emph{width $\to \infty$
first}, with $n$ and $L$ fixed (more precisely $L/N \to 0$ and
$n/N \to 0$); $1/\sqrt N$ output scaling with $O(1)$ learning rate ---
the $\gamma = 1/2$ endpoint of \cref{sec:gamma}, whereas $\gamma = 1$
makes $\Delta K = O(1)$ and feature learning survives; training time
$O(1)$, extendable to $t \to \infty$ when $\ntk(X,X)$ is positive
definite (the residual then decays exponentially and total parameter
motion stays bounded); for discrete-time gradient descent, additionally
$\eta < 2/\lambda_{\max}(\ntk(X,X))$ for stability (the constant
reappears as the edge of stability in \cref{sec:sgd}).

\emph{Derivation 2: effective-action power counting.} In the language of
\cref{sec:gp-defs}, the initialized network defines a nearly-Gaussian
theory: at $N = \infty$ the effective action for the outputs is quadratic
(\cref{sec:nngp-derivation}), and finite width adds interaction vertices
suppressed by powers of $1/N$ (\cref{sec:eft}). Now watch what sources
kernel motion. Differentiating \cref{eq:ntk-flow} once more,
\begin{equation}
  \label{eq:dntk}
  \partial_t K_t(\bv{x},\bv{x}')
  = -\sum_\mu \mathrm{d}K_t(\bv{x},\bv{x}',\bv{x}_\mu)\,
  \bigl(f_t(\bv{x}_\mu) - y_\mu\bigr),
  \quad
  \mathrm{d}K = \nabla_\theta\bigl(\nabla_\theta f(\bv{x})^{\top}
  \nabla_\theta f(\bv{x}')\bigr)^{\top}\nabla_\theta f ,
\end{equation}
the same hierarchy as \cref{eq:kernel-ode}, one level up. The
three-point object $\mathrm{d}K$ is, in the free theory, an odd-order
correlator of the Gaussian fields and their parameter-derivatives ---
and its expectation and fluctuations are generated only by the
interaction vertices of the effective action, so power counting assigns
it $O(1/N)$ (growing to $O(L/N)$ with depth; this is the ``dNTK'' of
\cite{roberts2022principles}, where the counting is done diagram by
diagram). At $N = \infty$ the hierarchy truncates at its first level:
$K$ is a deterministic ``c-number,'' $\mathrm{d}K \to 0$, and
\cref{eq:ntk-flow} closes. The two heuristics are really one: the
displacement argument counts the same powers of $1/N$ trajectory-wise
that the effective action counts correlator-wise. And the comparison
with \cref{sec:mup} becomes sharp: in mean-field scaling the analogous
third-order object $T_t$ in \cref{eq:kernel-ode} is $O(1)$ --- the
hierarchy does not truncate, and that non-truncation \emph{is} feature
learning.

\emph{Derivation 3: Bayesian learning.} Forget gradients entirely and
learn as a Bayesian: put the infinite-width prior over functions ---
which is exactly the NNGP of \cref{sec:nngp-derivation} --- and condition
on the $n$ observations (with Gaussian likelihood, noise variance
$\varsigma^2$). Conditioning a Gaussian measure yields a Gaussian
measure: the posterior is again a GP, with
\begin{equation}
  \label{eq:gp-posterior}
  \bar f(\bv{x}) = K(\bv{x},X)\bigl(K(X,X) + \varsigma^2 I\bigr)^{-1}
  \bv{y},
  \qquad
  K_{\mathrm{post}} = K - K(\cdot,X)\bigl(K(X,X) + \varsigma^2
  I\bigr)^{-1} K(X,\cdot),
\end{equation}
kernel ridge regression once more. Here nothing ever moves at all:
``training'' is linear algebra applied to a fixed kernel, because the
Gaussian family is closed under conditioning --- the Bayesian version of
``the features are frozen.'' The scaling regime is transparent in this
language too: exact Gaussianity of the prior needs $N \to \infty$ at
fixed $n$ and $L$, and at large-but-finite width the $n$ observations
perturb the per-neuron distributions at $O(n/N)$, the same ratio as
before. One honest subtlety: the Bayesian posterior mean uses the
\emph{NNGP} kernel $K^{(L)}$, while gradient-flow training uses the
\emph{NTK} --- two different kernels (related by the layer recursion
below), so infinite-width Bayesian inference and infinite-width gradient
training are two different kernel regressions with the same structural
story. Only for the last layer trained alone do they coincide.

For completeness, the NTK is computed by a layer recursion parallel to
\cref{eq:nngp-recursion}. Writing $\ntk^{(\ell)}$ for the NTK of the
network truncated at layer $\ell$ and
$\dot K^{(\ell+1)}(\bv{x},\bv{x}') = \sigma_w^2\,\E[\phi'(u)\phi'(v)]$
(expectation over the same bivariate Gaussian as in
\cref{eq:nngp-recursion}),
\begin{equation}
  \ntk^{(1)} = K^{(1)},
  \qquad
  \ntk^{(\ell+1)}(\bv{x},\bv{x}') = K^{(\ell+1)}(\bv{x},\bv{x}')
  + \dot K^{(\ell+1)}(\bv{x},\bv{x}')\;\ntk^{(\ell)}(\bv{x},\bv{x}') :
\end{equation}
each layer contributes its NNGP kernel, and everything upstream is
multiplied by a derivative factor. Note for later: at coincident inputs
this derivative factor is exactly the $\chi$ of
\cref{sec:criticality} --- the same number governs gradient propagation at
initialization and the depth-accumulation of the NTK, one reason
criticality and trainability are so tightly linked.

What the free theory buys: exact, closed-form predictions for training and
generalization (next subsection) that genuinely match wide networks. What it
misses: the features never move, so nothing resembling representation
learning occurs; predictions for tasks that \emph{require} features (see the
parity example below) are far too pessimistic. Interestingly, fine-tuning of
large language models appears to occur in a near-linearized regime
\cite{malladi2023kernel} --- the free theory is not only a toy.

\subsubsection{Eigenlearning: from architecture to inductive bias}
\label{sec:eigenlearning}

The NTK limit reduces deep learning to kernel regression. Kernel regression,
in turn, admits closed-form generalization estimates. This is the
centerpiece of the block, following \cite{simon2023eigenlearning} (this is
the example shown in Fig.~1b of the reading; the same result was first
obtained by replica methods in \cite{bordelon2020spectrum,canatar2021spectral}).

Set up the eigenproblem on the data distribution: with inputs
$\bv{x} \sim \pi$,
\begin{equation}
  \int K(\bv{x},\bv{x}')\,\phi_i(\bv{x}')\,\pi(d\bv{x}')
  = \lambda_i\,\phi_i(\bv{x}),
  \qquad
  K(\bv{x},\bv{x}') = \sum_i \lambda_i\,\phi_i(\bv{x})\,\phi_i(\bv{x}'),
\end{equation}
with $\{\phi_i\}$ orthonormal in $L^2(\pi)$ and eigenvalues sorted
decreasingly. Expand the target as $f^{*} = \sum_i v_i\,\phi_i$. Kernel
ridge regression (KRR) with ridge $\delta \ge 0$ on a training set of $n$
samples --- here and in all kernel sections, $n$ counts training samples,
playing the role $P$ played in \cref{sec:statmech} --- produces an estimate
$\hat f$; since KRR is linear in the target, its average action is
mode-by-mode multiplication, and we can define the \emph{learnability} of
mode $i$:
\begin{equation}
  L_i \;=\; \frac{\E\,\langle \hat f, \phi_i\rangle_{L^2(\pi)}}{v_i}
  \;\in\; [0, 1],
\end{equation}
the fraction of that mode's coefficient the learner captures on average.

\begin{theorem}[omniscient risk estimate \cite{simon2023eigenlearning}]
\label{thm:eigenlearning}
To leading order in large $n$ (and exactly in the appropriate scaling
limit),
\begin{equation}
  L_i = \frac{\lambda_i}{\lambda_i + \kappa},
  \qquad
  \text{where } \kappa \ge 0 \text{ solves}\quad
  \frac{\delta}{\kappa} + \sum_i \frac{\lambda_i}{\lambda_i + \kappa} = n,
\end{equation}
and the expected test risk on target $f^* = \sum_i v_i \phi_i$ with label
noise of variance $\varsigma^2$ is
\begin{equation}
  \mathcal{E} = E_0 \left[\, \sum_i (1 - L_i)^2\, v_i^2
  \;+\; \varsigma^2 \right],
  \qquad
  E_0 = \frac{n}{\,n - \sum_i L_i^2\,}.
\end{equation}
\end{theorem}

Unpack this slowly; every piece has physical content.
\begin{itemize}
  \item \emph{The conservation law.} Summing the first equation:
  $\sum_i L_i = n - \delta/\kappa \le n$, with equality at zero ridge.
  Training on $n$ samples buys \emph{exactly $n$ units of learnability},
  which the kernel allocates greedily to its top eigenmodes --- $L_i
  \approx 1$ for $\lambda_i \gg \kappa$, $L_i \approx 0$ for
  $\lambda_i \ll \kappa$. The constant $\kappa$ is the sharp threshold.
  \item \emph{Architecture $\Rightarrow$ generalization, quantitatively.}
  The causal chain is now explicit: the architecture fixes the NTK
  (\cref{def:ntk} and its layer recursion); the NTK and the data
  distribution fix the spectrum $\{\lambda_i, \phi_i\}$; the spectrum and
  $n$ fix the learnabilities; the learnabilities and the target fix the
  test risk. Every arrow is computable. This is the reading's claim that
  ``inductive bias follows from architecture through NTK eigenstructure,''
  with no hand-waving left.
  \item \emph{The overfitting factor.} $E_0 \geq 1$ multiplies
  \emph{everything}, including the noise floor. It diverges when
  $\sum_i L_i^2 \to n$, which happens at the interpolation threshold of a
  finite-rank kernel: this divergence \emph{is} double descent, derived
  rather than observed. In the language of \cref{sec:phase-transitions},
  the interpolation threshold is a phase transition in the ratio of
  samples to effective capacity, with $E_0$ the diverging
  susceptibility-like quantity --- smoothed at any finite size and ridge,
  singular in the limit; the lab measures its finite-size shadow directly.
  \item \emph{No replicas were harmed.} The modern derivation runs through
  the resolvent of the empirical kernel matrix: one shows that the
  test-risk formula depends on the data only through
  $\mathrm{tr}\,(K(X,X) + \delta I)^{-1}$, whose deterministic equivalent
  (a standard random-matrix statement) produces the self-consistency
  equation for $\kappa$ --- morally, a leave-one-out argument: adding one
  sample perturbs the resolvent by a rank-one update. The replica
  computation \cite{canatar2021spectral} gives the identical answer, and
  comparing the two derivations side by side is a methods lesson in
  itself.
\end{itemize}

\begin{remark}[the free Fermi gas]
\label{rem:fermi}
A \emph{free Fermi gas} is a collection of non-interacting particles
subject to an exclusion principle: each single-particle energy level can be
occupied by at most one particle, so at low temperature the particles
simply fill the levels from the bottom up to a sharp threshold (the
\emph{Fermi surface}), which smears slightly at positive temperature. The
structure $\{L_i \in [0,1]$, $\sum_i L_i = n$, $L_i =
\lambda_i/(\lambda_i + \kappa)\}$ is exactly this
\cite{simon2023eigenlearning}: learnabilities are occupation numbers
(between $0$ and $1$, by ``exclusion''), $-\log\lambda_i$ are the
single-particle energy levels, and $\kappa$ plays the role of the chemical
potential --- the Lagrange multiplier that enforces the particle number
$n$. Modes ``fill'' from the top of the spectrum, with a smeared Fermi
surface at $\lambda_i \approx \kappa$. For physicists this fixes the whole
picture in one line, and makes the scaling-law calculation of
\cref{sec:scaling} feel inevitable: power-law density of states
$\Rightarrow$ power-law observables.
\end{remark}

\subsubsection{An application with bite: parity}
\label{sec:parity}

Work on the Boolean cube, $\bv{x}$ uniform on $\{\pm 1\}^d$. Any
permutation-symmetric kernel (e.g.\ the NTK of a fully-connected network,
which depends only on $\bv{x}\cdot\bv{x}'/d$) is diagonalized by the
\emph{parities} $\chi_S(\bv{x}) = \prod_{i \in S} x_i$, with eigenvalues
$\lambda_{|S|}$ depending only on the degree $k = |S|$. Two facts do all the
work: eigenvalues decay rapidly with degree ($\lambda_k = O(d^{-k})$ for the
FCN NTK), and the degeneracy at degree $k$ is $\binom{d}{k}$. By
\cref{thm:eigenlearning}, learning a single $k$-subset parity requires
$\kappa \lesssim \lambda_k$, i.e.
\begin{equation}
  n \;\gtrsim\; d^{\,k} :
\end{equation}
kernels pay for high-order parities at a sample cost exponential in the
order. A network that \emph{learns features} --- concentrates its
first-layer weights on the $k$ relevant coordinates --- escapes this
scaling. This is the cleanest quantitative statement of \emph{why feature
learning matters}, and it sets the agenda for Block 2: find the scaling
limit in which features actually move. (We revisit exactly this example in
\cref{sec:mup}, where rich training visibly rebuilds the kernel around the
target parity.)

\begin{exercise}[NTK spectra and inductive bias]
\label{ex:ntk-circle}
\important
For a depth-$\ell$ ReLU fully-connected network, compute the NTK spectrum on
the unit circle (inputs $\bv{x} = (\cos\vartheta,\sin\vartheta)$;
eigenfunctions are Fourier modes by rotation invariance). Show how depth
reshapes the spectral decay, and hence the inductive bias, at fixed
parameter count.
\end{exercise}

\begin{hint}
Rotation invariance means the kernel is a function of the angle between
inputs alone; diagonalize by Fourier series in that angle, and use the
closed-form arc-cosine kernels (or \texttt{neural-tangents}
\cite{novak2020neural}) for the layer recursion.
\end{hint}

\begin{exercise}[Conservation of learnability]
\label{ex:conservation}
Derive $\sum_i L_i = n - \delta/\kappa$ from \cref{thm:eigenlearning}, and
show $0 \le L_i \le 1$. For a rank-$m$ kernel (only $m$ nonzero
eigenvalues) at $\delta \to 0$, show that $E_0$ diverges as $n \to m$, and
interpret.
\end{exercise}

%===============================================================================
\section{Lab: run the numbers}
\label{sec:lab}

Nothing in \cref{sec:eigenlearning} requires trust: every object in
\cref{thm:eigenlearning} is computable exactly in small synthetic settings,
and the theory's claims about \emph{real, finite} networks can be checked
directly. The goal of the lab is to reproduce the logic of Fig.~1b of
\cite{simon2026scientific} (equivalently Fig.~2 of
\cite{simon2023eigenlearning}) end to end: closed-form prediction first,
kernel regression second, trained networks third, all overlaid. This is the
single most persuasive ``physics-style theory works here'' experience the
day offers, so we run it hands-on rather than show plots.

\subsubsection*{Protocol}

\begin{enumerate}
  \item \textbf{Domain.} Discretized unit circle ($m$ equispaced points,
  $m \sim 128$; $m$ is unrelated to the teacher matrix $M$ of \cref{sec:dln}) and/or the Boolean cube $\{\pm 1\}^d$ with $d \approx 8$
  ($2^d = 256$ points). Both are small enough that the kernel eigensystem is
  exact by direct diagonalization of the $m \times m$ kernel matrix against
  the uniform measure --- no quadrature errors to argue about.
  \item \textbf{Kernel.} Exact NTK of a fully-connected architecture via
  \texttt{neural-tangents} \cite{novak2020neural}; sweep depth
  $\in \{1, 2, 4\}$ at fixed activation. Diagonalize; inspect the spectrum.
  On the circle, group eigenvalues by frequency; on the cube, by degree.
  \item \textbf{Theory.} Solve the self-consistency equation of
  \cref{thm:eigenlearning} for $\kappa(n)$ at each $n$ (a monotone scalar
  root-find; bisection is fine). Compute $L_i(n)$, $E_0(n)$, and the
  predicted test MSE for three targets: a single low-frequency mode, a
  single high-frequency/high-degree mode, and a mixture.
  \item \textbf{Experiment 1: kernel regression.} For each $n$: draw a
  training set, solve KRR exactly, measure test MSE on the remaining
  points; average over $\sim 50$ draws. Overlay on theory. Expected result:
  agreement at the percent level, including the double-descent spike where
  $E_0$ blows up.
  \item \textbf{Experiment 2: real networks.} Train finite-width
  fully-connected networks (width $\sim 500$, small learning rate, MSE
  loss, $1/\sqrt{N}$ output scaling to stay lazy) at several $n$; overlay.
  Agreement is strikingly good already at moderate width --- and the gap
  \emph{is} interesting: it previews finite-width corrections
  (\cref{sec:eft}).
  \item \textbf{Probes.} (i) Depth sweep: watch the spectrum, and with it
  the mode learnabilities, move --- \emph{architecture visibly steering
  inductive bias}. (ii) Push a parity target to higher degree at fixed $n$
  and watch learnability collapse as predicted in \cref{sec:parity}.
  (iii) Track $E_0(n)$ near the interpolation threshold and locate double
  descent before measuring it.
\end{enumerate}

\begin{callout}[warning]
Common failure modes: forgetting that eigenvalues must be computed against
the \emph{data measure} (normalize the kernel matrix by $1/m$); degenerate
eigenvalues on symmetric domains (group modes of equal frequency/degree
before comparing learnabilities, or the arbitrary basis within an eigenspace
scrambles per-mode plots); too few dataset draws at small $n$ (the variance
is genuinely large there --- that is what $E_0$ is telling you). All runs
take under a minute on CPU at these sizes.
\end{callout}

\begin{callout}[tip]
If time permits, keep the trained networks from Experiment 2: the lab
extension in \cref{sec:mup} re-trains them in the \emph{rich} regime on the
XOR target and measures the empirical NTK at checkpoints, closing the loop
between the two halves of the day.
\end{callout}

%===============================================================================
\section{Block 2: Insightful limits --- criticality, finite width, parameterizations}
\label{sec:block2}

The perceptron block used the oldest trick in statistical physics: take a
limit in which fluctuations die and a law of large numbers takes over. This
block is about the limits available to \emph{deep} networks --- in depth, in
width, and in the joint scaling of parameterization with width --- and about
what each limit reveals. The through-line: \emph{which} limit you take is
not bookkeeping; it decides whether the limiting theory learns features at
all.

\subsection{Signal propagation and the edge of chaos}
\label{sec:criticality}

Before training enters at all, there is a question of \emph{initialization}:
for a very deep network, does information about the input even reach the
output? This is a solvable dynamical-systems problem --- depth plays the
role of time --- and its answer produced one of the cleanest predictive
successes in the field.

\subsubsection{The variance and correlation maps}

Random fully-connected network on inputs $\bv{x} \in \R^d$, hidden width
$N$ (both as in \cref{sec:ntk}), weights
$\mathcal{N}(0, \sigma_w^2/N)$, biases $\mathcal{N}(0, \sigma_b^2)$,
activation $\phi$. At large width, the NNGP recursion
\cref{eq:nngp-recursion} becomes a deterministic dynamical system for the
statistics of preactivations across depth. For a single input, the
preactivation variance $q^\ell$ obeys
\begin{equation}
  \label{eq:variance-map}
  q^{\ell+1} = \mathcal{V}(q^{\ell})
  := \sigma_w^2\; \E_{z\sim\mathcal{N}(0,1)}
  \bigl[\phi(\sqrt{q^{\ell}}\,z)^2\bigr] + \sigma_b^2 ,
\end{equation}
which for saturating $\phi$ has a stable fixed point $q^{*}$ reached in a
few layers. For a \emph{pair} of inputs, the normalized correlation
$c^\ell$ of their preactivations obeys a second map
$c^{\ell+1} = \mathcal{C}(c^{\ell})$, computed from the bivariate Gaussian
expectation in \cref{eq:nngp-recursion} at variance $q^*$. The map
$\mathcal{C}$ always has a fixed point at $c = 1$ (identical inputs stay
identical); everything hinges on its \emph{stability}, governed by the slope
\begin{equation}
  \label{eq:chi}
  \chi \;=\; \mathcal{C}'(1)
  \;=\; \sigma_w^{2}\, \E_{z\sim\mathcal{N}(0,1)}
  \bigl[\phi'(\sqrt{q^{*}}\,z)^{2}\bigr].
\end{equation}
($\chi$ is also exactly the mean-square gradient amplification per layer ---
the forward and backward pictures are two faces of the same number.)

\subsubsection{Two phases and a critical line}

\begin{itemize}
  \item $\chi < 1$ (\emph{ordered}): $c = 1$ is stable. Any two inputs
  collapse toward perfect correlation exponentially in depth,
  $1 - c^\ell \sim \chi^{\ell}$: the deep network maps everything to nearly
  the same point. Gradients vanish at the same exponential rate.
  \item $\chi > 1$ (\emph{chaotic}): $c = 1$ is unstable. Nearby inputs
  decorrelate exponentially --- infinitesimally different inputs diverge, the
  signature of chaos. Gradients explode.
  \item $\chi = 1$ (\emph{criticality}): the correlation map is marginal,
  correlations decay only \emph{algebraically}, and the depth scale
  \begin{equation}
    \xi_c = -1/\log \chi
  \end{equation}
  diverges. Information about input geometry survives to arbitrary depth.
\end{itemize}
The condition $\chi(\sigma_w, \sigma_b) = 1$ defines the \emph{critical
line} in the $(\sigma_w^2, \sigma_b^2)$ plane --- the ``edge of chaos''
\cite{poole2016exponential,schoenholz2017deep}. In the taxonomy of
\cref{sec:phase-transitions}, this is a genuine \emph{continuous} phase
transition in a deep network: control parameters $(\sigma_w^2,
\sigma_b^2)$, order-parameter-like observable $c^*$ changing character
across the line, a characteristic scale $\xi_c \approx 1/(1 - \chi)$
diverging at criticality (exponent $1$), and exponential decay giving way
to algebraic decay exactly at the critical point --- the standard
signatures, here with \emph{depth} playing the role of distance.

\begin{callout}[tip]
The payoff, and it is the day's cleanest: \emph{trainable depth tracks
$\xi_c$}. Networks initialized off the critical line are untrainable beyond
depth $\sim \xi_c$, quantitatively, as a phase diagram in
$(\sigma_w^2, \sigma_b^2)$ against depth \cite{schoenholz2017deep}.
Initializing \emph{on} the line --- plus a stronger ``dynamical isometry''
condition on the input-output Jacobian spectrum, achieved with orthogonal
weights --- allowed \cite{xiao2018dynamical} to train vanilla CNNs
\emph{10{,}000 layers deep}, with no batch normalization and no skip
connections. The theory predicted the recipe before the experiment
confirmed it. This is what ``physics-style theory'' should mean: a phase
diagram computed from the ensemble, then an engineering feat read off it.
\end{callout}

\begin{exercise}[Edge of chaos for $\tanh$]
\label{ex:edge-of-chaos}
For $\phi = \tanh$: (a) show that \cref{eq:variance-map} has a unique
stable fixed point $q^*$ for any $\sigma_w, \sigma_b$; (b) compute
$\chi(\sigma_w, \sigma_b)$ and find the critical line (numerically beyond
$\sigma_b = 0$; show $\sigma_w = 1$, $q^* = 0$ is the critical point on the
axis); (c) linearize $\mathcal{C}$ near $c = 1$ and derive
$1 - c^{\ell} \sim \chi^{\ell}$, hence $\xi_c = -1/\log\chi$ and its
divergence at criticality. (d) Repeat for ReLU and explain what is
degenerate about it. \note{ReLU is positively homogeneous, so
$\mathcal{V}$ is linear: at $\sigma_b = 0$, $q^{\ell+1} =
(\sigma_w^2/2)\,q^\ell$, and criticality forces $\sigma_w^2 = 2$ --- He
initialization --- with no ordered phase to speak of, only marginal
behavior. Criticality is the reason that initialization is default, whether
or not its users know it.}
\end{exercise}

\subsection{Finite width as an effective field theory}
\label{sec:eft}

The infinite-width network is a free (Gaussian) theory: solvable, and
incapable of feature learning. Real networks are finite. The natural physics
move is perturbation theory in $1/N$ around the Gaussian limit, and it works
--- this is the program of Roberts, Yaida \& Hanin
\cite{roberts2022principles}. In this subsection (and only here) we follow
their convention and write $n$ for width, $L$ for depth.

\subsubsection{The simplest honest calculation}

Rather than survey, we compute. Take the simplest possible deep network: a
deep \emph{linear} network of width $n$, $z^{(\ell+1)} = W^{(\ell)}
z^{(\ell)}$, weights i.i.d.\ $\mathcal{N}(0, 1/n)$, single input, and track
the normalized squared activation
\begin{equation}
  q_\ell = \frac{1}{n}\,\bigl|z^{(\ell)}\bigr|^2 .
\end{equation}
Conditional on layer $\ell$, each component of $z^{(\ell+1)}$ is Gaussian
with variance $q_\ell$, so
\begin{equation}
  \E\bigl[q_{\ell+1} \,\big|\, q_\ell\bigr] = q_\ell,
  \qquad
  \E\bigl[q_{\ell+1}^2 \,\big|\, q_\ell\bigr]
  = q_\ell^2\left(1 + \frac{2}{n}\right),
\end{equation}
(the second moment of a normalized $\chi^2_n$). The mean is preserved, but
fluctuations \emph{compound multiplicatively}: iterating,
\begin{equation}
  \label{eq:qfluct}
  \frac{\E[q_L^2]}{\E[q_L]^2} = \left(1 + \frac{2}{n}\right)^{L}
  \;\approx\; e^{\,2L/n}.
\end{equation}
Two lines of algebra, one structural conclusion: the size of
non-Gaussian fluctuations after $L$ layers is controlled not by $1/n$ or by
$L$ separately but by the \emph{aspect ratio}
\begin{equation}
  \frac{L}{n} \;=\; \frac{\text{depth}}{\text{width}} ,
\end{equation}
which is therefore the \emph{effective coupling} of the theory.

\subsubsection{The general picture}

The full RYH analysis promotes this to nonlinear networks and to the
tangent kernel. Preactivation distributions at finite width are
\emph{nearly Gaussian}: the \emph{connected four-point function} --- the
fourth cumulant, i.e.\ the part of a four-point correlation not accounted
for by products of two-point functions, which vanishes identically for
Gaussian distributions and is therefore the natural first measure of
non-Gaussianity --- is $O(1/n)$, and it obeys a depth recursion whose
solution grows linearly in $L$ --- so again $L/n$ is the expansion
parameter, now for the full interacting theory.
The regimes:
\begin{itemize}
  \item $L/n \to 0$: free theory; the kernel limit of \cref{sec:ntk} with
  Gaussian fluctuations. Solvable, no features.
  \item $L/n$ small but finite: interacting but perturbative. The NTK now
  \emph{fluctuates} across initializations at $O(\sqrt{L/n})$, and ---
  crucially --- \emph{moves during training} at $O(L/n)$: the leading
  perturbative trace of feature learning. This is the regime where
  effective theory is both nontrivial and controlled, and empirically where
  well-tuned architectures live.
  \item $L/n \sim 1$: strongly coupled; perturbation theory fails, and
  experimentally such networks train poorly --- large depth-to-width ratios
  are avoided in practice, an observation the expansion explains.
\end{itemize}
The rigorous counterpart for shallow networks is the CLT of
\cite[\S20.4]{spiliopoulos2025foundations}: the empirical distribution of
neurons is its mean-field limit plus $N^{-1/2}$ Gaussian fluctuations. And
the same structure can be organized as a dictionary between networks and
interacting field theories, with finite-width non-Gaussianities playing the
role of interaction couplings \cite{halverson2021neural} --- an alternative
language worth a pointer for QFT-minded students.

\begin{exercise}[Deep linear fluctuations]
\label{ex:deep-linear-fluct}
Derive both conditional moments above (use rotational invariance of $W$),
verify \cref{eq:qfluct}, and compute the connected fourth moment of a
single output component at depth $L$ to leading order in $1/n$. Where does
the calculation break down as $L/n \to 1$?
\end{exercise}

\subsection{The mean-field limit: lazy vs.\ rich, and the road to $\mu$P}
\label{sec:mup}

Primary text: \cite[ch.~20]{spiliopoulos2025foundations} (compare their
ch.~19 for the NTK regime). We work out the one-hidden-layer heuristics in
full and cite the rigorous statements; the book is the students' bridge to
the proofs.

\subsubsection{The $N\to\infty$ mean-field limit, heuristically}

Everything in \cref{sec:ntk} used the $1/\sqrt{N}$ output normalization.
Change one exponent. The one-hidden-layer network in \emph{mean-field
normalization} is
\begin{equation}
  f(\bv{x};\theta) = \frac{1}{N}\sum_{j=1}^{N} c_j\,
  \sigma(\bv{w}_j\cdot\bv{x})
  = \int \rho(\theta;\bv{x})\,\mu^N(d\theta),
  \qquad
  \rho\bigl((c,\bv{w});\bv{x}\bigr) = c\,\sigma(\bv{w}\cdot\bv{x}),
\end{equation}
where $\mu^N = \frac1N \sum_j \delta_{\theta_j}$ is the \emph{empirical
measure} of the neurons $\theta_j = (c_j, \bv{w}_j)$, $j = 1, \dots, N$ (we index neurons by $j$; $n$ remains the sample count). This rewriting is the
whole idea: the network is a linear functional of $\mu^N$, so the population
loss
\begin{equation}
  \mathcal{L}(\mu) = \tfrac12 \int \Delta_\mu(\bv{x})^2\, \pi(d\bv{x}),
  \qquad \Delta_\mu = f_\mu - f^{*},
\end{equation}
is a smooth (quadratic) functional on the space of measures, and the
$N \to \infty$ limit is a law of large numbers for interacting particles.

Now run SGD with an $O(1)$ learning rate $\eta$ for $\lfloor N t \rfloor$
steps (hydrodynamic time: one unit of $t$ costs $N$ steps). The gradient of
the loss with respect to one neuron carries a $1/N$ from the normalization,
which combines with the $N$ steps per unit time to produce an $O(1)$
displacement: averaging over the data, each neuron drifts as
\begin{equation}
  \dot\theta_j = -\eta\, \nabla_\theta \Psi(\theta_j; \mu^N_t),
  \qquad
  \Psi(\theta;\mu) = \int \Delta_\mu(\bv{x})\,
  \rho(\theta;\bv{x})\,\pi(d\bv{x}) .
\end{equation}
Read this as physics: each particle moves in a \emph{single-particle
potential} $\Psi$ whose shape is determined by all the other particles only
through the measure $\mu$ --- the defining structure of a \emph{mean-field}
system (``mean field'': each unit interacts with the rest only through
their collective, averaged effect, never pairwise with individuals). The
continuity equation --- the transport PDE stating that a density changes
only by probability mass flowing in or out under the velocity field, here
$-\nabla_\theta\Psi$ --- then gives, in the limit,
\begin{equation}
  \label{eq:mft}
  \partial_t \mu_t = \nabla_\theta \cdot
  \bigl(\mu_t\, \nabla_\theta \Psi(\theta;\mu_t)\bigr),
\end{equation}
which is precisely the \emph{Wasserstein gradient flow} of $\mathcal{L}$
--- steepest descent on the space of probability measures, with distances
measured in the optimal-transport (Wasserstein) metric, where the cost of
moving from one measure to another is the cheapest way to transport its
mass. Along the flow,
\begin{equation}
  \frac{d}{dt}\mathcal{L}(\mu_t)
  = -\int \bigl|\nabla_\theta \Psi(\theta;\mu_t)\bigr|^2\,\mu_t(d\theta)
  \;\le\; 0 .
\end{equation}
Both SGD noise and initialization randomness vanish in the limit --- the
dynamics are deterministic, with the $1/N$ factor playing the role usually
played by learning-rate decay. Rigorous versions:
\cite[Thm.~20.4, Cor.~20.5]{spiliopoulos2025foundations} (tightness $+$
martingale decomposition $+$ uniqueness by contraction), propagation of
chaos in \cite{sirignano2020mean}, and the parallel routes of
\cite{mei2018mean,chizat2018global,rotskoff2022trainability}.

\subsubsection{Kernel evolution: how the moving kernel encodes features}
\label{sec:kernel-evolution}

By the chain rule, the network output \emph{always} obeys an equation of
NTK form:
\begin{equation}
  \label{eq:output-ode}
  \partial_t f_t(\bv{x}) = -\int A_t(\bv{x},\bv{x}')\,
  \Delta_t(\bv{x}')\,\pi(d\bv{x}'),
  \qquad
  A_t(\bv{x},\bv{x}') = \int \nabla_\theta\rho(\theta;\bv{x})\cdot
  \nabla_\theta\rho(\theta;\bv{x}')\,\mu_t(d\theta),
\end{equation}
explicitly
$A_t(\bv{x},\bv{x}') = \int \bigl[\sigma(\bv{w}\cdot\bv{x})
\sigma(\bv{w}\cdot\bv{x}') + c^2\,\sigma'(\bv{w}\cdot\bv{x})
\sigma'(\bv{w}\cdot\bv{x}')\,\bv{x}\cdot\bv{x}'\bigr]\,\mu_t(d\theta)$.
The \emph{only} difference from the NTK regime is that there the kernel is
frozen, $A_t = A_0$, while here it moves with $\mu_t$
\cite[Rem.~20.1]{spiliopoulos2025foundations}. Lazy and rich training solve
the same output equation with a static versus a live kernel; that is the
entire dichotomy in one sentence.

How does the kernel move? Differentiate $A_t$ along \cref{eq:mft}
(integrate by parts against the test function
$\nabla\rho(\bv{x})\cdot\nabla\rho(\bv{x}')$):
\begin{equation}
  \label{eq:kernel-ode}
  \partial_t A_t(\bv{x},\bv{x}') = -\int T_t(\bv{x},\bv{x}',\bv{x}'')\,
  \Delta_t(\bv{x}'')\,\pi(d\bv{x}''),
\end{equation}
where
\begin{equation}
  T_t(\bv{x},\bv{x}',\bv{x}'') = \int \nabla_\theta\bigl(
  \nabla_\theta\rho(\bv{x})\cdot\nabla_\theta\rho(\bv{x}')\bigr)
  \cdot \nabla_\theta\rho(\bv{x}'')\,\mu_t(d\theta).
\end{equation}
The kernel's motion is \emph{sourced by the residual} $\Delta_t$, through a
third-order object $T_t$; and $T_t$ evolves in turn via a fourth-order
kernel, and so on --- a hierarchy that does not close, of BBGKY type (the
situation familiar from kinetic theory, where the evolution of the
$k$-particle correlation function always involves the $(k{+}1)$-particle
one, and some truncation or closure must be chosen).
\emph{Feature learning is precisely the motion of $A_t$.} The hierarchy's
non-closure is not a defect but a signpost: DMFT (\cref{sec:dmft}) is the
systematic bookkeeping that closes it self-consistently in deep networks,
and the RYH $1/n$ expansion (\cref{sec:eft}) is its leading perturbative
truncation.

\begin{example}[silent alignment; kernels learn the task's geometry]
\label{ex:silent-alignment}
During rich training, the top eigenspace of $A_t$ rotates toward the
target: kernel--task alignment rises \emph{early}, before the loss falls
appreciably, and for small initialization the trained network is well
approximated by kernel regression with the \emph{final}, aligned kernel
\cite{atanasov2022silent}. Concrete instance on the lab's domain: XOR,
$f^{*}(\bv{x}) = x_i x_j$ on the Boolean cube. At initialization $A_0$ is
permutation-symmetric and the parity mode is nearly unlearnable at moderate
$n$ (\cref{sec:parity}); under rich training, first-layer weights
concentrate on coordinates $i, j$, and the top eigenfunction of $A_t$
becomes $\approx x_i x_j$. Feature learning is visible as \emph{spectral
reorganization of the kernel}: the network builds itself the kernel that
makes the task easy.
\end{example}

\begin{callout}[tip]
Lab extension (30 min if time permits): re-train the lab's networks in the
rich regime ($1/N$ scaling, or a small output multiplier) on XOR; measure
the empirical NTK at checkpoints; feed its spectrum into
\cref{thm:eigenlearning}; watch the predicted learnability of the target
mode grow. This ties the day's two halves together: eigenlearning predicts
generalization \emph{given} a kernel, feature learning explains where good
kernels come from.
\end{callout}

\begin{exercise}[Kernel evolution]
\label{ex:kernel-evolution}
Derive \cref{eq:output-ode} and \cref{eq:kernel-ode} from \cref{eq:mft}.
Exhibit the fourth-order kernel governing $\partial_t T_t$ and explain why
the hierarchy does not close. Where does the NTK regime's frozen kernel
reappear in this language?
\end{exercise}

\subsubsection{Between the regimes: $\gamma$-scalings collapse to the kernel limit}
\label{sec:gamma}

The NTK normalization is $N^{-1/2}$; mean-field is $N^{-1}$. What about
everything in between? Scale the output as $N^{-\gamma}$ with learning rate
$\eta/N^{2(1-\gamma)}$, $\gamma \in [1/2, 1)$. The result
\cite[\S20.6]{spiliopoulos2025foundations} is a collapse: \emph{every} such
$\gamma$ has the same infinite-width limit --- the NTK limit. The refinement
is an asymptotic expansion of the output,
\begin{equation}
  h_t^{N,\gamma} \approx h_t + \sum_{j\ge1} N^{-j(1-\gamma)}\, Q_t^{j}
  + N^{-(\gamma - 1/2)}\, e^{-At}\,\mathcal{G} + \cdots,
\end{equation}
with $h_t$ the common NTK-limit trajectory, $Q^j$ deterministic
corrections, and $\mathcal{G}$ a Gaussian: the leading \emph{fluctuation}
variance scales as $N^{-2(\gamma - 1/2)}$, decreasing monotonically as
$\gamma \to 1$. Test accuracy in the book's MNIST/CIFAR experiments improves
monotonically in $\gamma$ accordingly. Moral: the kernel regime is an entire
\emph{basin} of scalings, and mean-field ($\gamma = 1$) is the distinguished
endpoint --- the unique scaling with $O(1)$ kernel motion, and minimal
fluctuations besides. This is the sharp, one-parameter version of the
lazy/rich dichotomy of \cite{chizat2019lazy} and of Fig.~2 of the reading,
where the same finite network is pushed between regimes by an output
multiplier.

\subsubsection{Deeper networks: forced learning-rate scalings, and $\mu$P}
\label{sec:mup-deep}

Repeat the exercise for two hidden layers ($N_1$, then $N_2$) in mean-field
scaling \cite[\S20.5]{spiliopoulos2025foundations}, and something new
happens: a nontrivial training limit \emph{forces layer-dependent learning
rates},
\begin{equation}
  \eta_C = \frac{N_2}{N_1}, \qquad \eta_{W^1} = 1, \qquad \eta_{W^2} = N_2 .
\end{equation}
With equal learning rates across layers, the network provably \emph{does not
train} in the limit --- some layers freeze or diverge. Under mild additional
conditions the limit trains to a \emph{global} minimum. The book remarks
that the analysis ``suggests a closed form formula for the choice of the
learning rate hyperparameters,'' and that the formula, more than the limit
itself, is the practically valuable output. That sentence, generalized, is
$\mu$P.

\begin{callout}[note]
Dictionary slide. Mean-field scaling $=$ $\mu$P restricted to one hidden
layer. The $abc$-parameterizations of \cite{yang2021tensor} scale
multipliers ($a$), initializations ($b$), and learning rates ($c$) as
powers of width for \emph{every} layer of a network of \emph{any} depth,
and classify all limits: each choice flows to a kernel limit, diverges, or
--- on one distinguished manifold --- retains $O(1)$ feature learning in
every layer at once. That distinguished choice is the \emph{maximal-update
parameterization} ($\mu$P); the collapse result of \cref{sec:gamma} is its
one-parameter shadow, and the forced learning rates above are its depth-2
instance. The book's ``closed-form formula for learning rates'' is $\mu$P's
table of width exponents.
\end{callout}

\begin{callout}[tip]
Payoff: hyperparameter transfer ($\mu$Transfer). In $\mu$P, optimal
learning rates and related hyperparameters become width-independent, so one
tunes a small proxy model and transfers to the production-scale model
\cite{yang2022tensor} --- used in real large-model training. This closes the
loop on the reading's \S2.4, and is the alignment-adjacent moral of the
block: \emph{principled small-to-large extrapolation is possible exactly
when you know the right scaling limit.}
\end{callout}

%===============================================================================
\section{Block 3: Dynamics and macroscopic laws}
\label{sec:block3}

Block 2 identified the limits; this block is about what actually happens
during training in them --- and about the macroscopic empirical laws
(scaling laws, edge of stability) that any candidate theory must explain.
The reading's \S2.3 makes the methodological point: the most lawful
quantities in deep learning are \emph{aggregate} statistics, and physics has
always built theories from such laws (Kepler before Newton). We take three
case studies in ascending order of empirical contact.

\subsection{Dynamical mean-field theory of feature learning}
\label{sec:dmft}

The kernel hierarchy of \cref{ex:kernel-evolution} did not close. For deep
networks at large width there is a systematic closure: the dynamical
mean-field theory of \cite{bordelon2022self}. We teach the \emph{structure}
--- what the order parameters are and what the theory predicts --- and leave
the derivation as flagged optional reading (it uses the Martin--Siggia--Rose
formalism, a path-integral representation of stochastic dynamics in which
correlation functions are computed from a generating functional and its
saddle point; students who know the Keldysh formalism from nonequilibrium
quantum theory, or dynamical treatments of disordered systems, will
recognize the machinery).

\begin{itemize}
  \item \emph{Order parameters.} For each layer $\ell$, the feature kernel
  $\Phi^{\ell}(t, s)$ and gradient kernel $G^{\ell}(t, s)$: inner products
  of representations (respectively backpropagated signals) at training
  times $t$ and $s$. These are two-time correlation functions, as in the
  dynamics of glasses and other aging systems (cf.\ the glassy-phase
  discussion of \cref{sec:replica-method}): training breaks
  time-translation invariance --- the system never forgets when it started
  --- so correlations depend on both times, not just their difference, and
  one clock is not enough.
  \item \emph{Single-site picture.} At the DMFT saddle point, neurons
  decouple: each hidden unit becomes an i.i.d.\ stochastic process driven
  by Gaussian fields whose correlators are the kernels, and the kernels are
  determined self-consistently as averages over the single-site process.
  This is the same self-consistency logic as \cref{thm:eigenlearning}'s
  $\kappa$, upgraded from a scalar to a set of two-time functions.
  \item \emph{Lazy--rich interpolation.} A richness parameter $\gamma_0$
  (the output multiplier, suitably scaled --- related in spirit to, but
  distinct from, the width-scaling exponent $\gamma$ of \cref{sec:gamma})
  interpolates between frozen kernels ($\gamma_0 \to 0$: the NTK regime, kernels static in $t$) and
  fully rich dynamics ($\gamma_0 = O(1)$: kernels evolve, features form).
  The shallow mean-field limit of \cref{sec:mup} is the rigorous special
  case; DMFT extends it to depth and closes the kernel hierarchy
  self-consistently.
  \item \emph{What it predicts.} Learning curves and kernel evolution that
  match finite networks quantitatively at moderate width; and a controlled
  language for ``what changed inside the network during training'' --- the
  physics counterpart of interpretability's questions, and the source of
  candidate \emph{monitoring observables} (kernels as order parameters
  whose motion signals structure formation).
\end{itemize}

\subsection{SGD as a stochastic process; the edge of stability}
\label{sec:sgd}

\subsubsection{The SDE picture and its honest success}

Minibatch SGD with learning rate $\eta$ and batch size $B$ is, between
steps, a noisy difference equation; at small $\eta$ it is well modeled by
the SDE
\begin{equation}
  d\theta = -\nabla L(\theta)\, dt
  + \sqrt{\frac{\eta}{B}}\; \Sigma^{1/2}(\theta)\, dW,
\end{equation}
where $\Sigma$ is the covariance of per-sample gradients. A caution before
any equilibrium intuition kicks in: \emph{Langevin dynamics} --- gradient
descent plus \emph{isotropic} white noise of fixed temperature $T$ --- has
the special property that its stationary probability measure is the Gibbs
measure $\propto e^{-L/T}$, so its long-time behavior is equilibrium
statistical mechanics. SGD noise is anisotropic and state-dependent
($\Sigma$ varies over parameter space and vanishes at interpolating
minima), so no such closed-form stationary measure exists, and equilibrium
reasoning does not apply; resist the temptation. The scaling consequence is
immediate and useful: the drift and the noise depend on $(\eta, B)$ only
through the ratio $\eta/B$, so doubling both leaves the trajectory
approximately invariant --- the \emph{linear scaling rule}, used in practice
to parallelize training, with $\sqrt{B}$ variants for adaptive optimizers.
A modest, honest success: a one-line stochastic-process argument organizing
a family of expensive experiments.

\subsubsection{The anomaly: progressive sharpening and the edge of stability}

For \emph{full-batch} gradient descent the SDE picture degenerates to
gradient flow --- and here an experimental anomaly forced better theory.
Descent on a quadratic with curvature $\lambda$ --- here and below,
$\lambda$ is a \emph{Hessian} eigenvalue, not a kernel eigenvalue as in
\cref{sec:eigenlearning} --- obeys
$\theta_{t+1} = (1 - \eta\lambda)\,\theta_t$, stable iff
$\lambda < 2/\eta$. Naively, networks should sit far from this bound.
Instead \cite{cohen2021gradient} found, robustly across architectures:
the top Hessian eigenvalue (``sharpness'') \emph{rises} during training
(progressive sharpening) until it hits $2/\eta$, then \emph{stays pinned
there} while the loss continues to decrease --- non-monotonically, through
sustained oscillation. Gradient descent typically operates \emph{at} the
edge of stability, in flat contradiction of the gradient-flow picture and
of every classical convergence proof's hypotheses.

\subsubsection{The resolution: central flows}

The dynamics at the edge separate into fast oscillation along the top
Hessian direction and slow drift transverse to it --- a textbook multiscale
structure. Averaging over the fast oscillation yields a \emph{central flow}
\cite{cohen2025understanding}: the time-averaged trajectory follows
\begin{equation}
  \dot\theta = -\nabla L(\theta) \;-\; c(t)\,\nabla S(\theta),
\end{equation}
gradient flow on the loss plus an \emph{implicit penalty on the sharpness}
$S(\theta)$, whose coefficient $c(t)$ (proportional to the oscillation
variance) is determined self-consistently by the requirement that $S$ stay
pinned at $2/\eta$. The oscillations are not a nuisance around the ``real''
dynamics; they \emph{are} the mechanism that regulates curvature. In the
reading's summary: we now have a mathematical understanding of what the
learning rate does in full-batch gradient descent --- it sets the curvature
budget, and the dynamics self-organize to saturate it. For physicists: slow
manifold plus fast oscillation, averaging theory, and a self-stabilized
critical state.

\begin{exercise}[Edge of stability warm-up]
\label{ex:eos}
(a) Derive the $2/\eta$ threshold for GD on a quadratic. (b) On a cubic
$L(\theta) = \tfrac{\lambda}{2}\theta_1^2 + \tfrac{g}{2}\theta_1^2\theta_2
+ \dots$ with $\lambda$ near $2/\eta$, show that oscillation in $\theta_1$
produces a systematic drift in $\theta_2$ proportional to the oscillation
variance, in the direction that \emph{reduces} $\lambda$ if
$\partial\lambda/\partial\theta_2 = g > 0$. This is the germ of the central
flow's sharpness regulation.
\end{exercise}

\subsection{Scaling laws from solvable models}
\label{sec:scaling}

The empirical facts \cite{kaplan2020scaling,hoffmann2022training}: over many
orders of magnitude, test loss falls as a power law in model size $N$, data
$D$, and compute $C$; compute-optimal training balances $N$ and $D$ along a
predictable frontier. These laws underwrite the entire economics of large
models --- and the current scientific basis of \emph{capability forecasting},
which is why an alignment course cares. The reading is admirably blunt:
``Why does test loss decay as a power law\dots\ We still do not know!'' But
theory gets remarkably far, and the machinery is exactly what we built in
\cref{sec:eigenlearning}.

\subsubsection{Data scaling from the eigenlearning fixed point}

Suppose the kernel spectrum on the data decays as a power law,
$\lambda_i \sim i^{-\alpha}$ ($\alpha > 1$) --- empirically true for natural
data across modalities, with $\alpha$ measurable. (This $\alpha$ is the
conventional symbol for the spectral-decay exponent; it has nothing to do
with the load $\alpha = P/N$ of Block~1, which is not in play here.) Solve the
self-consistency equation of \cref{thm:eigenlearning} at zero ridge:
\begin{equation}
  n = \sum_i \frac{\lambda_i}{\lambda_i + \kappa}
  = \sum_i \frac{1}{1 + \kappa\, i^{\alpha}}
  \;\approx\; \int_0^\infty \frac{di}{1 + \kappa\, i^{\alpha}}
  \;=\; C_\alpha\, \kappa^{-1/\alpha},
\end{equation}
so
\begin{equation}
  \kappa(n) \;\sim\; n^{-\alpha}:
\end{equation}
the Fermi level sweeps down the spectrum as a power of the dataset size, and
the number of learned modes is $\sim n$. Now take a target whose
coefficients also decay as a power law, $v_i^2 \sim i^{-\beta}$
($\beta > 1$). The risk splits at the Fermi surface $i_* \sim n$:
\begin{equation}
  \mathcal{E}(n) \;\approx\;
  \underbrace{\sum_{i \lesssim n} \bigl(\kappa\, i^{\alpha}\bigr)^2 i^{-\beta}}_{\text{learned modes, residual error}}
  \;+\;
  \underbrace{\sum_{i \gtrsim n} i^{-\beta}}_{\text{unlearned tail}}
  \;\sim\; n^{-(\beta - 1)},
\end{equation}
both pieces contributing at the same order (a short computation; do it).
Special case worth memorizing: a target drawn generically from the kernel's
own prior has $v_i^2 \propto \lambda_i$, giving
\begin{equation}
  \label{eq:data-scaling}
  \mathcal{E}(n) \;\sim\; n^{-(\alpha - 1)} .
\end{equation}
\emph{Derived, not fitted}: the learning-curve exponent is a computable
functional of the data spectrum and target alignment (in the kernel
literature, ``capacity and source conditions''
\cite{bordelon2020spectrum,cui2023error}). Since students built $\kappa(n)$
numerically in the lab, this derivation is a corollary, not a new
machine.

\subsubsection{Joint scaling and the compute-optimal frontier}

Model size enters by truncating the spectrum: a model with $P$ parameters
can resolve roughly $P$ modes. The random-matrix model of
\cite{maloney2022solvable} makes this exact --- power-law features, finite
model rank, full RMT treatment --- and yields a loss surface
$\mathcal{L}(N, D)$ with power laws in both arguments, a
parameters$\leftrightarrow$data duality, and a derived compute-optimal
frontier: loss is minimized at fixed compute by scaling data and parameters
together (as Chinchilla found empirically \cite{hoffmann2022training}).
A complementary organizing scheme \cite{bahri2021explaining} distinguishes
\emph{resolution-limited} regimes (exponents set by spectral decay, as
above) from \emph{variance-limited} ones (exponents $= 1$, set by
fluctuation averaging). And the DMFT of \cref{sec:dmft} extends the whole
story into the rich regime and into \emph{time}: a dynamical model of
scaling laws \cite{bordelon2024dynamical} yields the loss as a joint power
law in training time, width, and data, including compute-optimal tradeoffs
--- theory catching up to the quantities practitioners actually budget.
What theory does not yet deliver: exponents for \emph{individual
capabilities} rather than aggregate loss, and a first-principles account of
why natural data is power-law in the first place --- both are live research
problems, and both are exactly where the phase-transitions day picks up.

\begin{callout}[note]
Alignment payoff, stated plainly: scaling-law theory is the current
scientific basis of capability forecasting. Its known holes ---
per-capability exponents, emergent discontinuities hiding under smooth
aggregate loss --- are not footnotes; they are the precise technical form of
the question ``will the next model surprise us?''
\end{callout}

\begin{exercise}[Scaling exponents]
\label{ex:scaling}
Carry out both sums in the risk split above and verify that they are of the
same order, $n^{-(\beta-1)}$. Then redo the calculation at finite ridge
$\delta > 0$ and with label noise $\varsigma^2 > 0$, and show that noise
introduces a variance-limited term $\sim \varsigma^2\, n^{-1}\cdot
n\,/\,(n - \sum_i L_i^2)$ --- interpret its interplay with the
resolution-limited term.
\end{exercise}

%===============================================================================
\section{Closing: universality and alignment synthesis}
\label{sec:closing}

This closing session is a discussion, not a lecture, built around the
reading's \S2.5 and the day's accumulated evidence.

\subsection{Universality}

Three empirical claims, in increasing order of ambition. First,
\emph{universal inductive biases}: at matched compute, architectures as
different as ConvNets and vision transformers reach closely matched
performance, and different diffusion architectures generate nearly identical
images from the same noise seed --- details of the microscopic model wash
out. Second, \emph{universal structure in data}: the no-free-lunch theorem
says no learner beats another averaged over all tasks, so deep learning's
success must be a fact about \emph{our} data --- and natural data does show
recurring structure: power-law spectra (the empirical input to
\cref{eq:data-scaling}), Zipf's law in text, hierarchical compositionality
across modalities. Third, and boldest, \emph{universal representations}:
networks trained across different initializations, architectures,
objectives, and even modalities learn increasingly similar internal
representations as they scale --- the ``Platonic representation hypothesis''
\cite{huh2024platonic} --- though the strength of the claim depends on the
similarity metric, a caveat students should press on.

For physicists the resonance is with the renormalization group ---
systematic coarse-graining, under which microscopic details wash out and
the large-scale behavior of a system is governed by a small number of
``relevant'' parameters, so that microscopically different systems flow to
the same effective description: behavior shared across microscopically
different systems should depend only on what those systems share, and
therefore admit a description simpler than any one of them. That is the deepest reason to expect a mechanics of learning to
exist at all. And it carries a curious corollary the reading dwells on: if
internal structure primarily reflects the data rather than the model, then
in studying neural networks we are, in the end, studying the structure of
our data --- including, for language models, the structure of us.

\subsection{The Discretization Hypothesis, revisited}

Recall the claim from Session 0: practical networks are noisy, finite
approximations of infinite-size models, so finite-size effects typically
cost performance rather than enable it. The day supplied the limits the
hypothesis refers to: the kernel limit (\cref{sec:ntk}), criticality at
infinite depth-scale (\cref{sec:criticality}), the mean-field limit and
$\mu$P (\cref{sec:mup}), DMFT (\cref{sec:dmft}). Test the morning's
candidate falsifiers against them. Two productive pressure points: SGD
noise (is its regularization benefit achievable in a limit? --- the
$\gamma$-expansion of \cref{sec:gamma} says fluctuations only hurt, but
that is a statement about one family of limits), and capacity constraints
(a finite model that \emph{cannot} memorize might generalize where an
infinite one interpolates noise --- or the infinite model at the right
ridge does better still; \cref{ex:scaling} is relevant). The hypothesis
survives the day, in our experience, but not without a fight, and the fight
is the point.

\subsection{The alignment through-line}

Assemble the day's payoffs into one argument:
\begin{enumerate}
  \item \emph{Small-to-large extrapolation works when the right limit is
  known.} Trainability phase diagrams predicted 10{,}000-layer training;
  $\mu$Transfer tunes frontier models on small proxies. These are
  existence proofs that scaling behavior is \emph{lawful}, not folklore.
  \item \emph{Capability forecasting has a scientific core}
  (\cref{sec:scaling}) --- derived exponents from measurable spectral data
  --- and known holes: per-capability exponents, and discontinuous
  capability emergence beneath smooth aggregate loss.
  \item \emph{Monitoring needs order parameters.} The day produced several
  observables that move when structure forms inside a network: kernel
  spectra and kernel--task alignment (\cref{ex:silent-alignment}), DMFT's
  two-time kernels (\cref{sec:dmft}), sharpness pinned at $2/\eta$
  (\cref{sec:sgd}). The phase-transitions and SLT days sharpen this idea
  into a program: find observables that move \emph{before} the capability
  arrives.
\end{enumerate}
The gap between items 2 and 3 --- forecasting what emerges, and detecting it
as it forms --- is where the next days of this course live.

\subsection{Further reading}

For Block 1: \cite{engel2001statistical} remains the canonical textbook on
the statistical mechanics of learning, with the perceptron calculations of
\cref{sec:statmech} in full; \cite{zdeborova2016statistical} is the modern
review connecting to inference and rigorous results;
\cite{bahri2020statistical} surveys the statistical-mechanics view of deep
learning specifically and makes a good bridge from this day to the current
literature. For the kernel side, \cite{simon2023eigenlearning} is short and
readable, and \cite{canatar2021spectral} gives the replica route with
applications to real datasets. For Block 2:
\cite{roberts2022principles} is the effective-theory book (chapters 4--5
for the four-point recursion behind \cref{sec:eft});
\cite[chs.~19--20]{spiliopoulos2025foundations} is the rigorous
mathematical treatment of both regimes, at exactly the level of these
notes' heuristics; \cite{yang2022tensor} documents $\mu$Transfer in
practice. For Block 3: \cite{bordelon2022self} for DMFT,
\cite{cohen2025understanding} for central flows (unusually well-written),
and \cite{maloney2022solvable,bordelon2024dynamical} for scaling-law
theory in the static and dynamical settings respectively.

%===============================================================================
\section{Appendix: the replica calculation, guided}
\label{sec:replica-appendix}

This appendix expands the outline of \cref{sec:replica-method} into a
step-by-step scaffold for the full quenched calculation (result:
$\varepsilon \simeq 0.625/\alpha$, \cref{eq:quenched-eps}), at the level of
structure rather than every Gaussian integral; the full details are in
\cite[chs.~2--3]{engel2001statistical}. It is intended as a guided
afternoon for students who want the real thing, and overlaps deliberately
with the main text so it can be worked standalone.

\begin{enumerate}
  \item \emph{Replicate.} For integer $n$, write $\E Z^n = \E \int
  \prod_{a=1}^{n} d\nu(\bv{w}^a) \prod_{\mu, a}
  \Theta(y^\mu\, \bv{w}^a \cdot \bv{x}^\mu)$: an $n$-fold copy of the
  system coupled only through the shared data.
  \item \emph{Average over data.} Each example couples to the replicas only
  through the $n+1$ jointly Gaussian fields $u^a = \bv{w}^a\cdot\bv{x}/
  \sqrt N$ and $v = \bv{w}^*\cdot\bv{x}/\sqrt N$, whose covariance is
  fixed by the \emph{overlap matrix}
  \begin{equation}
    q^{ab} = \frac{\bv{w}^a \cdot \bv{w}^b}{N},
    \qquad
    R^a = \frac{\bv{w}^a \cdot \bv{w}^*}{N}.
  \end{equation}
  The data average factorizes across the $P$ examples, each contributing
  the same function of $(q^{ab}, R^a)$: the disorder is gone, at the price
  of coupling the replicas.
  \item \emph{Change variables to the order parameters.} Insert delta
  functions fixing $q^{ab}$ and $R^a$; the $\bv{w}$-integrals become an
  entropic factor $e^{N S(q, R)}$, and $\E Z^n \asymp \int dq\, dR\;
  e^{N[\,S(q,R) + \alpha\, \mathcal{A}(q,R)\,]}$: a saddle-point problem
  over matrices, the direct generalization of \cref{eq:annealed-saddle}.
  \item \emph{Replica symmetry.} Take the ansatz $q^{ab} = q$ for all
  $a \neq b$, $R^a = R$. The free entropy becomes an explicit function of
  two scalars (involving the Gaussian tail function
  $H(x) = \int_x^\infty \frac{dz}{\sqrt{2\pi}} e^{-z^2/2}$), analytically
  continued in $n$, and differentiated at $n = 0$.
  \item \emph{Use the symmetry of Gibbs learning.} A student drawn
  uniformly from the version space is statistically indistinguishable from
  the teacher (both are ``typical'' consistent vectors), which forces
  $q = R$ at the saddle: the overlap between two Gibbs students equals
  their overlap with the teacher. One scalar remains.
  \item \emph{Solve and expand.} The saddle-point equation for $R(\alpha)$
  is transcendental; solve numerically for the full learning curve, and
  expand at large $\alpha$ to extract $\varepsilon \simeq 0.625/\alpha$
  --- same exponent as the annealed result \cref{eq:annealed-eps},
  corrected constant. Compare the two curves; note that the annealed
  approximation is exact as $\alpha \to 0$ and never off by more than the
  constant at large $\alpha$: for this problem, annealing is a good
  approximation, and knowing \emph{when} that is true (realizable rules,
  no glassiness) is part of the education.
\end{enumerate}

Where replica symmetry breaks (unrealizable rules, binary weights at high
load, and most glassy optimization problems), the ansatz in step 4 must be
enriched; \cite{zdeborova2016statistical} explains both the physics and the
rigorous status of the resulting hierarchy.

%===============================================================================
\section*{Selected solutions}

\begin{solution}[ex:annealed]
The wedge computation: $(\bv{w}\cdot\bv{x},\bv{w}^*\cdot\bv{x})/\sqrt{N}$
is a centered Gaussian pair with unit variances and correlation $R$; the
probability of opposite signs is the fraction of the plane between the two
hyperplanes' normal directions, $\arccos(R)/\pi$. The slice entropy: the
sphere $|\bv{w}|^2 = N$ at fixed $\bv{w}\cdot\bv{w}^* = NR$ is a sphere of
radius $\sqrt{N(1-R^2)}$ in $N-1$ dimensions, whose log-volume per site is
$\tfrac12\log(1-R^2) + \text{const}$. Laplace's method on
$\int dR\, e^{N[s(R) + \alpha\log(1-\varepsilon(R))]}$ gives
\cref{eq:annealed-saddle}. Near $R = 1$, with $R = \cos\pi\varepsilon$:
$G \approx \log(\pi\varepsilon) - \alpha\varepsilon$, maximized at
$\varepsilon = 1/\alpha$. At small $\alpha$, expand instead around $R = 0$:
$G \approx -R^2/2 + \alpha\log(1/2 + R/\pi + \dots)$, giving
$R \approx 2\alpha/\pi \propto \alpha$; the learning curve interpolates
smoothly between the two regimes.
\end{solution}

\begin{solution}[ex:ising]
(a) $\log\binom{N}{pN} = N[-p\log p - (1-p)\log(1-p)] + O(\log N)$ by
Stirling; set $p = (1+R)/2$.
(b) With $\delta = 1 - R$: $H\bigl(1 - \tfrac{\delta}{2}\bigr) =
\tfrac{\delta}{2}\log\tfrac{2}{\delta} + \tfrac{\delta}{2} + O(\delta^2)$,
while $\arccos(1 - \delta) = \sqrt{2\delta}\,(1 + O(\delta))$, so
$G_{\mathrm{I}}(1 - \delta) = -\tfrac{\alpha\sqrt{2\delta}}{\pi}
+ O(\delta\log(1/\delta)) < 0$ for small $\delta$ and any $\alpha > 0$:
$G_{\mathrm{I}}$ dips below $G_{\mathrm{I}}(1) = 0$ in a neighborhood of
$R = 1$, which is therefore a local maximum.
(c) Numerics (bisection on $\alpha$ for $\max_R G_{\mathrm{I}} = 0$):
$\alpha_c^{\mathrm{ann}} \approx 1.448$, with $R^* \approx 0.698$,
$\varepsilon^* \approx 0.254$; the slope of $\Phi_{\mathrm{ann}}$ jumps
from $\log(1 - \varepsilon^*) \approx -0.293$ to $0$; the interior local
maximum persists to a spinodal at $\alpha \approx 1.73$ and then
disappears. (Quenched: $\alpha_c \approx 1.245$
\cite{gyorgyi1990first}.)
(d) Implementation notes: propose single-spin flips; reject any flip that
violates a constraint; measure $R$ (hence $\varepsilon(R)$) along the
chain. Expect strong hysteresis near and beyond $\alpha_c$: runs
initialized at random $R \approx 0$ stay near the interior branch well
past $1.245$, while runs initialized near the teacher stay at $R = 1$ ---
the two-basin structure made visible.
\end{solution}

\begin{solution}[ex:edge-of-chaos]
(a) For $\tanh$, $\mathcal{V}$ is increasing and concave in $q^\ell$ with
$\mathcal{V}(0) = \sigma_b^2$ and $\mathcal{V}$ bounded by $\sigma_w^2 +
\sigma_b^2$, so it has a unique fixed point, reached monotonically.
(b) $\chi = \sigma_w^2 \E[\mathrm{sech}^4(\sqrt{q^*} z)]$; on the axis
$\sigma_b = 0$: for $\sigma_w \le 1$ the only fixed point is $q^* = 0$
where $\chi = \sigma_w^2$, so criticality is at $\sigma_w = 1$; for
$\sigma_w > 1$, $q^* > 0$ and $\chi$ crosses $1$ along a curve
$\sigma_b^2(\sigma_w^2)$ found numerically.
(c) Writing $c^\ell = 1 - \epsilon_\ell$: $\epsilon_{\ell+1} =
\chi\,\epsilon_\ell + O(\epsilon_\ell^2)$, so $\epsilon_\ell \sim
\chi^\ell = e^{-\ell/\xi_c}$ with $\xi_c = -1/\log\chi \to \infty$ as
$\chi \to 1$. At $\chi = 1$ the quadratic term controls the map and decay
is algebraic, $\epsilon_\ell \sim 1/\ell$.
(d) For ReLU, $\phi(\lambda u) = \lambda\phi(u)$ makes $\mathcal{V}$
linear: $q^{\ell+1} = (\sigma_w^2/2)\, q^\ell + \sigma_b^2$. There is no
nontrivial fixed-point structure to tune: at $\sigma_b = 0$ variance is
exactly preserved iff $\sigma_w^2 = 2$ (He initialization), where the
network sits at marginality automatically --- ReLU networks do not have an
ordered/chaotic transition to cross so much as a razor's edge to balance
on, and the correlation map approaches $c = 1$ algebraically.
\end{solution}

\begin{solution}[ex:conservation]
Summing $L_i = \lambda_i/(\lambda_i+\kappa)$ over $i$ and using the
constraint gives $\sum_i L_i = n - \delta/\kappa$ directly. Bounds:
$\lambda_i, \kappa \ge 0$. For a rank-$m$ kernel at $\delta \to 0$: if
$n < m$, the constraint forces $\kappa > 0$ and each $L_i < 1$; as
$n \to m^{-}$, $\kappa \to 0$, every $L_i \to 1$, so
$\sum_i L_i^2 \to m \to n$ and $E_0$ diverges. Interpretation: at the
interpolation threshold the model uses every sample to pin every mode, no
averaging remains, and the estimator's variance explodes --- double descent.
Past the threshold ($n > m$) the excess samples restore averaging and
$E_0$ falls back toward $1$.
\end{solution}

\begin{solution}[ex:deep-linear-fluct]
Condition on $z^{(\ell)}$ and use rotational invariance to set
$z^{(\ell)} = |z^{(\ell)}|\,e_1$: then $z^{(\ell+1)}_i = |z^{(\ell)}|\,
W_{i1}$ are i.i.d.\ $\mathcal{N}(0, q_\ell)$, so $q_{\ell+1} =
q_\ell\cdot\chi^2_n/n$ with $\E[\chi^2_n/n] = 1$,
$\E[(\chi^2_n/n)^2] = 1 + 2/n$. Independence of the $W^{(\ell)}$ across
layers lets the second moments multiply, giving \cref{eq:qfluct}. For a
single output component $z^{(L)}_1$: conditionally Gaussian with variance
$q_{L-1}$, so $\E[(z_1^{(L)})^4] = 3\,\E[q_{L-1}^2] =
3(1+2/n)^{L-1}\E[q_0^2]$, and the excess kurtosis is
$3[(1+2/n)^{L-1} - 1] \approx 6L/n$ at leading order: the connected
four-point function grows linearly in depth over width. The expansion
parameter is $L/n$; when $L/n \sim 1$ the neglected higher cumulants are
$O(1)$ and the perturbative (nearly-Gaussian) description fails.
\end{solution}

\begin{solution}[ex:kernel-evolution]
For any functional $F[\mu] = \int g(\theta)\,\mu(d\theta)$,
\cref{eq:mft} and integration by parts give $\dot F = -\int
\nabla g\cdot\nabla\Psi\, d\mu_t$. Take $g = \rho(\theta;\bv{x})$: since
$\nabla\Psi(\theta) = \int \Delta_t(\bv{x}')\,\nabla\rho(\theta;\bv{x}')\,
\pi(d\bv{x}')$, this is \cref{eq:output-ode} with the stated $A_t$. Take
$g = \nabla\rho(\bv{x})\cdot\nabla\rho(\bv{x}')$: this is
\cref{eq:kernel-ode} with the stated $T_t$. Iterating, the time derivative
of the $k$-point kernel involves the $(k+1)$-point kernel: each
differentiation inserts one more $\nabla\rho$, so no finite set of moments
closes --- the hierarchy is BBGKY-like, with $\mu_t$ itself the only closed
description. The NTK regime corresponds to the degenerate case where
$\mu_t = \mu_0$ to leading order (the $1/\sqrt N$ scaling makes
$\nabla\Psi$ act on a scale where individual particles barely move):
then $A_t = A_0$ and the hierarchy truncates at its first level.
\end{solution}

\begin{solution}[ex:scaling]
Learned modes: $\sum_{i \lesssim n} \kappa^2 i^{2\alpha - \beta}$ with
$\kappa \sim n^{-\alpha}$; if $2\alpha - \beta > -1$ the sum is dominated
by its upper end, $\sim n^{-2\alpha}\, n^{2\alpha - \beta + 1} =
n^{1-\beta}$. Unlearned tail: $\sum_{i \gtrsim n} i^{-\beta} \sim
n^{1-\beta}$. Both $\sim n^{-(\beta - 1)}$. With noise: the term
$E_0\,\varsigma^2$ never decays below $\varsigma^2$; using $\sum_i L_i^2
\sim c\,n$ (compute $c < 1$ for a power-law spectrum), $E_0 = O(1)$, and
the risk floors out at $\sim \varsigma^2$ plus a variance-limited
correction. The observable consequence: learning curves cross over from the
resolution-limited power law $n^{-(\beta-1)}$ to a noise plateau at
$n_* \sim \varsigma^{-2/(\beta-1)}$; with explicit regularization the ridge
$\delta$ shifts $\kappa$ and trades bias against the noise term ---
optimizing it recovers the classical bias--variance-optimal rates of the
kernel literature \cite{cui2023error}.
\end{solution}

\begin{solution}[ex:eos]
(a) $\theta_{t+1} = (1-\eta\lambda)\theta_t$ gives $|1-\eta\lambda| < 1
\iff 0 < \lambda < 2/\eta$. (b) With $\theta_1$ oscillating at amplitude
$a$ about $0$ near the stability boundary, average the $\theta_2$ update:
$\E[\partial_{\theta_2} L] = \tfrac{g}{2}\E[\theta_1^2] =
\tfrac{g}{2}a^2$, so $\theta_2$ drifts at rate $-\eta g a^2/2$: downhill
in the direction that reduces the effective curvature
$\lambda_{\mathrm{eff}} = \lambda + g\theta_2$ when $g > 0$. Larger
oscillations produce faster sharpness reduction; the self-consistent state
pins $\lambda_{\mathrm{eff}}$ at $2/\eta$ with just enough oscillation
variance to hold it there --- the central flow's mechanism in miniature.
\end{solution}

\bibliographystyle{plain}
\bibliography{biblo}

\end{document}
